\section{未锚定账本事件的叙事补充：western-jin}

\label{sec:anchor-batch-two-western-jin}

以下事件均为此前正文尚未设置导航目标的账本记录。每一锚点置于来源、制度与地方条件的叙事中，避免把年表、政权分类或战报修辞当作社会全貌。

\subsection{制度与地方资源}

\eventanchor{JH-LEDGER-004-1}{永平元年（291年）·“贾后联络楚王司马玮等发动政变，杨骏被杀，杨氏外戚集团被清除”的前期动员与资源准备}账本记录西晋在永平元年（291年）发生““贾后联络楚王司马玮等发动政变，杨骏被杀，杨氏外戚集团被清除”的前期动员与资源准备”。条目把背景概括为当时的权力、资源与信息约束推动相关过程展开，并指出这一过程为“贾后联络楚王司马玮等发动政变，杨骏被杀，杨氏外戚集团被清除”提供可观察的条件。这个节点进入区域社会时，要经由文书、道路、军队、市场、寺院和家庭等不同中介；因此，政治行动的宣布、地方接收和日常生活的改变不能被视作同一时刻完成。

本条依据《晋书·惠帝纪》；《晋书·贾后传》；《晋书·杨骏传》；《资治通鉴·晋纪》，证据提示为“核心事项已有既有账本来源；本分解节点的参与者、执行范围和时间先后仍须逐案核验”。这些材料可支持年序、官方决定、战事或特定地点的线索，却难以直接统计所有人的财富、迁徙、信仰和动机。更稳妥的读法是区分诏令与实施、军报与人口、行政称谓与自我认同，并让地方材料的缺环限制解释的范围。

由此可见，该事件不是孤立的年表点，而是资源、信息与权力关系重新配置的一环。不同地区的官员、社群和家庭可能以不同方式承担征发、保护、迁移或宗教活动；现存来源只能照亮其中一部分，不能替沉默者制造统一的经验。

\eventanchor{JH-LEDGER-004-2}{永平元年（291年）·“贾后联络楚王司马玮等发动政变，杨骏被杀，杨氏外戚集团被清除”的命令、联盟或地方响应}账本记录西晋在永平元年（291年）发生““贾后联络楚王司马玮等发动政变，杨骏被杀，杨氏外戚集团被清除”的命令、联盟或地方响应”。条目把背景概括为当时的权力、资源与信息约束推动相关过程展开，并指出这一过程为“贾后联络楚王司马玮等发动政变，杨骏被杀，杨氏外戚集团被清除”提供可观察的条件。这个节点进入区域社会时，要经由文书、道路、军队、市场、寺院和家庭等不同中介；因此，政治行动的宣布、地方接收和日常生活的改变不能被视作同一时刻完成。

本条依据《晋书·惠帝纪》；《晋书·贾后传》；《晋书·杨骏传》；《资治通鉴·晋纪》，证据提示为“核心事项已有既有账本来源；本分解节点的参与者、执行范围和时间先后仍须逐案核验”。这些材料可支持年序、官方决定、战事或特定地点的线索，却难以直接统计所有人的财富、迁徙、信仰和动机。更稳妥的读法是区分诏令与实施、军报与人口、行政称谓与自我认同，并让地方材料的缺环限制解释的范围。

由此可见，该事件不是孤立的年表点，而是资源、信息与权力关系重新配置的一环。不同地区的官员、社群和家庭可能以不同方式承担征发、保护、迁移或宗教活动；现存来源只能照亮其中一部分，不能替沉默者制造统一的经验。

\eventanchor{JH-LEDGER-004-3}{永平元年（291年）·贾后联络楚王司马玮等发动政变，杨骏被杀，杨氏外戚集团被清除}账本记录西晋在永平元年（291年）发生“贾后联络楚王司马玮等发动政变，杨骏被杀，杨氏外戚集团被清除”。条目把背景概括为“贾后联络楚王司马玮等发动政变，杨骏被杀，杨氏外戚集团被清除”发生的政治与区域条件共同作用，并指出该节点改变后续资源、联盟或制度处置的条件。这个节点进入区域社会时，要经由文书、道路、军队、市场、寺院和家庭等不同中介；因此，政治行动的宣布、地方接收和日常生活的改变不能被视作同一时刻完成。

本条依据《晋书·惠帝纪》；《晋书·贾后传》；《晋书·杨骏传》；《资治通鉴·晋纪》，证据提示为“核心事项已有既有账本来源；本分解节点的参与者、执行范围和时间先后仍须逐案核验”。这些材料可支持年序、官方决定、战事或特定地点的线索，却难以直接统计所有人的财富、迁徙、信仰和动机。更稳妥的读法是区分诏令与实施、军报与人口、行政称谓与自我认同，并让地方材料的缺环限制解释的范围。

由此可见，该事件不是孤立的年表点，而是资源、信息与权力关系重新配置的一环。不同地区的官员、社群和家庭可能以不同方式承担征发、保护、迁移或宗教活动；现存来源只能照亮其中一部分，不能替沉默者制造统一的经验。

\eventanchor{JH-LEDGER-004-4}{永平元年（291年）·“贾后联络楚王司马玮等发动政变，杨骏被杀，杨氏外戚集团被清除”的执行中的空间与人员处置}账本记录西晋在永平元年（291年）发生““贾后联络楚王司马玮等发动政变，杨骏被杀，杨氏外戚集团被清除”的执行中的空间与人员处置”。条目把背景概括为“贾后联络楚王司马玮等发动政变，杨骏被杀，杨氏外戚集团被清除”发生的政治与区域条件共同作用，并指出该节点改变后续资源、联盟或制度处置的条件。这个节点进入区域社会时，要经由文书、道路、军队、市场、寺院和家庭等不同中介；因此，政治行动的宣布、地方接收和日常生活的改变不能被视作同一时刻完成。

本条依据《晋书·惠帝纪》；《晋书·贾后传》；《晋书·杨骏传》；《资治通鉴·晋纪》，证据提示为“核心事项已有既有账本来源；本分解节点的参与者、执行范围和时间先后仍须逐案核验”。这些材料可支持年序、官方决定、战事或特定地点的线索，却难以直接统计所有人的财富、迁徙、信仰和动机。更稳妥的读法是区分诏令与实施、军报与人口、行政称谓与自我认同，并让地方材料的缺环限制解释的范围。

由此可见，该事件不是孤立的年表点，而是资源、信息与权力关系重新配置的一环。不同地区的官员、社群和家庭可能以不同方式承担征发、保护、迁移或宗教活动；现存来源只能照亮其中一部分，不能替沉默者制造统一的经验。

\eventanchor{JH-LEDGER-004-5}{永平元年（291年）·“贾后联络楚王司马玮等发动政变，杨骏被杀，杨氏外戚集团被清除”的后续制度或军政调整}账本记录西晋在永平元年（291年）发生““贾后联络楚王司马玮等发动政变，杨骏被杀，杨氏外戚集团被清除”的后续制度或军政调整”。条目把背景概括为核心节点发生后仍须处理其遗留问题，并指出后续影响须在不同地区和材料类型中分别核验。这个节点进入区域社会时，要经由文书、道路、军队、市场、寺院和家庭等不同中介；因此，政治行动的宣布、地方接收和日常生活的改变不能被视作同一时刻完成。

本条依据《晋书·惠帝纪》；《晋书·贾后传》；《晋书·杨骏传》；《资治通鉴·晋纪》，证据提示为“核心事项已有既有账本来源；本分解节点的参与者、执行范围和时间先后仍须逐案核验”。这些材料可支持年序、官方决定、战事或特定地点的线索，却难以直接统计所有人的财富、迁徙、信仰和动机。更稳妥的读法是区分诏令与实施、军报与人口、行政称谓与自我认同，并让地方材料的缺环限制解释的范围。

由此可见，该事件不是孤立的年表点，而是资源、信息与权力关系重新配置的一环。不同地区的官员、社群和家庭可能以不同方式承担征发、保护、迁移或宗教活动；现存来源只能照亮其中一部分，不能替沉默者制造统一的经验。

\eventanchor{JH-LEDGER-004-6}{永平元年（291年）·“贾后联络楚王司马玮等发动政变，杨骏被杀，杨氏外戚集团被清除”的异源材料与影响范围的复核}账本记录西晋在永平元年（291年）发生““贾后联络楚王司马玮等发动政变，杨骏被杀，杨氏外戚集团被清除”的异源材料与影响范围的复核”。条目把背景概括为核心节点发生后仍须处理其遗留问题，并指出后续影响须在不同地区和材料类型中分别核验。这个节点进入区域社会时，要经由文书、道路、军队、市场、寺院和家庭等不同中介；因此，政治行动的宣布、地方接收和日常生活的改变不能被视作同一时刻完成。

本条依据《晋书·惠帝纪》；《晋书·贾后传》；《晋书·杨骏传》；《资治通鉴·晋纪》，证据提示为“核心事项已有既有账本来源；本分解节点的参与者、执行范围和时间先后仍须逐案核验”。这些材料可支持年序、官方决定、战事或特定地点的线索，却难以直接统计所有人的财富、迁徙、信仰和动机。更稳妥的读法是区分诏令与实施、军报与人口、行政称谓与自我认同，并让地方材料的缺环限制解释的范围。

由此可见，该事件不是孤立的年表点，而是资源、信息与权力关系重新配置的一环。不同地区的官员、社群和家庭可能以不同方式承担征发、保护、迁移或宗教活动；现存来源只能照亮其中一部分，不能替沉默者制造统一的经验。

\eventanchor{JH-LEDGER-005-1}{元康元年（291年）·“汝南王司马亮与卫瓘辅政后又被贾后借楚王司马玮诛杀，司马玮旋即被处死”的前期动员与资源准备}账本记录西晋在元康元年（291年）发生““汝南王司马亮与卫瓘辅政后又被贾后借楚王司马玮诛杀，司马玮旋即被处死”的前期动员与资源准备”。条目把背景概括为当时的权力、资源与信息约束推动相关过程展开，并指出这一过程为“汝南王司马亮与卫瓘辅政后又被贾后借楚王司马玮诛杀，司马玮旋即被处死”提供可观察的条件。这个节点进入区域社会时，要经由文书、道路、军队、市场、寺院和家庭等不同中介；因此，政治行动的宣布、地方接收和日常生活的改变不能被视作同一时刻完成。

本条依据《晋书·惠帝纪》；《晋书·汝南文成王亮传》；《晋书·卫瓘传》；《晋书·楚隐王玮传》，证据提示为“核心事项已有既有账本来源；本分解节点的参与者、执行范围和时间先后仍须逐案核验”。这些材料可支持年序、官方决定、战事或特定地点的线索，却难以直接统计所有人的财富、迁徙、信仰和动机。更稳妥的读法是区分诏令与实施、军报与人口、行政称谓与自我认同，并让地方材料的缺环限制解释的范围。

由此可见，该事件不是孤立的年表点，而是资源、信息与权力关系重新配置的一环。不同地区的官员、社群和家庭可能以不同方式承担征发、保护、迁移或宗教活动；现存来源只能照亮其中一部分，不能替沉默者制造统一的经验。

\eventanchor{JH-LEDGER-005-2}{元康元年（291年）·“汝南王司马亮与卫瓘辅政后又被贾后借楚王司马玮诛杀，司马玮旋即被处死”的命令、联盟或地方响应}账本记录西晋在元康元年（291年）发生““汝南王司马亮与卫瓘辅政后又被贾后借楚王司马玮诛杀，司马玮旋即被处死”的命令、联盟或地方响应”。条目把背景概括为当时的权力、资源与信息约束推动相关过程展开，并指出这一过程为“汝南王司马亮与卫瓘辅政后又被贾后借楚王司马玮诛杀，司马玮旋即被处死”提供可观察的条件。这个节点进入区域社会时，要经由文书、道路、军队、市场、寺院和家庭等不同中介；因此，政治行动的宣布、地方接收和日常生活的改变不能被视作同一时刻完成。

本条依据《晋书·惠帝纪》；《晋书·汝南文成王亮传》；《晋书·卫瓘传》；《晋书·楚隐王玮传》，证据提示为“核心事项已有既有账本来源；本分解节点的参与者、执行范围和时间先后仍须逐案核验”。这些材料可支持年序、官方决定、战事或特定地点的线索，却难以直接统计所有人的财富、迁徙、信仰和动机。更稳妥的读法是区分诏令与实施、军报与人口、行政称谓与自我认同，并让地方材料的缺环限制解释的范围。

由此可见，该事件不是孤立的年表点，而是资源、信息与权力关系重新配置的一环。不同地区的官员、社群和家庭可能以不同方式承担征发、保护、迁移或宗教活动；现存来源只能照亮其中一部分，不能替沉默者制造统一的经验。

\eventanchor{JH-LEDGER-005-3}{元康元年（291年）·汝南王司马亮与卫瓘辅政后又被贾后借楚王司马玮诛杀，司马玮旋即被处死}账本记录西晋在元康元年（291年）发生“汝南王司马亮与卫瓘辅政后又被贾后借楚王司马玮诛杀，司马玮旋即被处死”。条目把背景概括为“汝南王司马亮与卫瓘辅政后又被贾后借楚王司马玮诛杀，司马玮旋即被处死”发生的政治与区域条件共同作用，并指出该节点改变后续资源、联盟或制度处置的条件。这个节点进入区域社会时，要经由文书、道路、军队、市场、寺院和家庭等不同中介；因此，政治行动的宣布、地方接收和日常生活的改变不能被视作同一时刻完成。

本条依据《晋书·惠帝纪》；《晋书·汝南文成王亮传》；《晋书·卫瓘传》；《晋书·楚隐王玮传》，证据提示为“核心事项已有既有账本来源；本分解节点的参与者、执行范围和时间先后仍须逐案核验”。这些材料可支持年序、官方决定、战事或特定地点的线索，却难以直接统计所有人的财富、迁徙、信仰和动机。更稳妥的读法是区分诏令与实施、军报与人口、行政称谓与自我认同，并让地方材料的缺环限制解释的范围。

由此可见，该事件不是孤立的年表点，而是资源、信息与权力关系重新配置的一环。不同地区的官员、社群和家庭可能以不同方式承担征发、保护、迁移或宗教活动；现存来源只能照亮其中一部分，不能替沉默者制造统一的经验。

\eventanchor{JH-LEDGER-005-4}{元康元年（291年）·“汝南王司马亮与卫瓘辅政后又被贾后借楚王司马玮诛杀，司马玮旋即被处死”的执行中的空间与人员处置}账本记录西晋在元康元年（291年）发生““汝南王司马亮与卫瓘辅政后又被贾后借楚王司马玮诛杀，司马玮旋即被处死”的执行中的空间与人员处置”。条目把背景概括为“汝南王司马亮与卫瓘辅政后又被贾后借楚王司马玮诛杀，司马玮旋即被处死”发生的政治与区域条件共同作用，并指出该节点改变后续资源、联盟或制度处置的条件。这个节点进入区域社会时，要经由文书、道路、军队、市场、寺院和家庭等不同中介；因此，政治行动的宣布、地方接收和日常生活的改变不能被视作同一时刻完成。

本条依据《晋书·惠帝纪》；《晋书·汝南文成王亮传》；《晋书·卫瓘传》；《晋书·楚隐王玮传》，证据提示为“核心事项已有既有账本来源；本分解节点的参与者、执行范围和时间先后仍须逐案核验”。这些材料可支持年序、官方决定、战事或特定地点的线索，却难以直接统计所有人的财富、迁徙、信仰和动机。更稳妥的读法是区分诏令与实施、军报与人口、行政称谓与自我认同，并让地方材料的缺环限制解释的范围。

由此可见，该事件不是孤立的年表点，而是资源、信息与权力关系重新配置的一环。不同地区的官员、社群和家庭可能以不同方式承担征发、保护、迁移或宗教活动；现存来源只能照亮其中一部分，不能替沉默者制造统一的经验。

\eventanchor{JH-LEDGER-005-5}{元康元年（291年）·“汝南王司马亮与卫瓘辅政后又被贾后借楚王司马玮诛杀，司马玮旋即被处死”的后续制度或军政调整}账本记录西晋在元康元年（291年）发生““汝南王司马亮与卫瓘辅政后又被贾后借楚王司马玮诛杀，司马玮旋即被处死”的后续制度或军政调整”。条目把背景概括为核心节点发生后仍须处理其遗留问题，并指出后续影响须在不同地区和材料类型中分别核验。这个节点进入区域社会时，要经由文书、道路、军队、市场、寺院和家庭等不同中介；因此，政治行动的宣布、地方接收和日常生活的改变不能被视作同一时刻完成。

本条依据《晋书·惠帝纪》；《晋书·汝南文成王亮传》；《晋书·卫瓘传》；《晋书·楚隐王玮传》，证据提示为“核心事项已有既有账本来源；本分解节点的参与者、执行范围和时间先后仍须逐案核验”。这些材料可支持年序、官方决定、战事或特定地点的线索，却难以直接统计所有人的财富、迁徙、信仰和动机。更稳妥的读法是区分诏令与实施、军报与人口、行政称谓与自我认同，并让地方材料的缺环限制解释的范围。

由此可见，该事件不是孤立的年表点，而是资源、信息与权力关系重新配置的一环。不同地区的官员、社群和家庭可能以不同方式承担征发、保护、迁移或宗教活动；现存来源只能照亮其中一部分，不能替沉默者制造统一的经验。

\eventanchor{JH-LEDGER-005-6}{元康元年（291年）·“汝南王司马亮与卫瓘辅政后又被贾后借楚王司马玮诛杀，司马玮旋即被处死”的异源材料与影响范围的复核}账本记录西晋在元康元年（291年）发生““汝南王司马亮与卫瓘辅政后又被贾后借楚王司马玮诛杀，司马玮旋即被处死”的异源材料与影响范围的复核”。条目把背景概括为核心节点发生后仍须处理其遗留问题，并指出后续影响须在不同地区和材料类型中分别核验。这个节点进入区域社会时，要经由文书、道路、军队、市场、寺院和家庭等不同中介；因此，政治行动的宣布、地方接收和日常生活的改变不能被视作同一时刻完成。

本条依据《晋书·惠帝纪》；《晋书·汝南文成王亮传》；《晋书·卫瓘传》；《晋书·楚隐王玮传》，证据提示为“核心事项已有既有账本来源；本分解节点的参与者、执行范围和时间先后仍须逐案核验”。这些材料可支持年序、官方决定、战事或特定地点的线索，却难以直接统计所有人的财富、迁徙、信仰和动机。更稳妥的读法是区分诏令与实施、军报与人口、行政称谓与自我认同，并让地方材料的缺环限制解释的范围。

由此可见，该事件不是孤立的年表点，而是资源、信息与权力关系重新配置的一环。不同地区的官员、社群和家庭可能以不同方式承担征发、保护、迁移或宗教活动；现存来源只能照亮其中一部分，不能替沉默者制造统一的经验。

\eventanchor{JH-LEDGER-006-1}{元康六年至九年（296年至299年）·“氐族齐万年在关中起兵，西晋多次征讨并长期调动关陇军粮，至299年叛乱被平定”的前期动员与资源准备}账本记录西晋与关陇氐羌集团在元康六年至九年（296年至299年）发生““氐族齐万年在关中起兵，西晋多次征讨并长期调动关陇军粮，至299年叛乱被平定”的前期动员与资源准备”。条目把背景概括为当时的权力、资源与信息约束推动相关过程展开，并指出这一过程为“氐族齐万年在关中起兵，西晋多次征讨并长期调动关陇军粮，至299年叛乱被平定”提供可观察的条件。这个节点进入区域社会时，要经由文书、道路、军队、市场、寺院和家庭等不同中介；因此，政治行动的宣布、地方接收和日常生活的改变不能被视作同一时刻完成。

本条依据《晋书·惠帝纪》；《晋书·张方传》；《晋书·江统传》；《资治通鉴·晋纪》，证据提示为“核心事项已有既有账本来源；本分解节点的参与者、执行范围和时间先后仍须逐案核验”。这些材料可支持年序、官方决定、战事或特定地点的线索，却难以直接统计所有人的财富、迁徙、信仰和动机。更稳妥的读法是区分诏令与实施、军报与人口、行政称谓与自我认同，并让地方材料的缺环限制解释的范围。

由此可见，该事件不是孤立的年表点，而是资源、信息与权力关系重新配置的一环。不同地区的官员、社群和家庭可能以不同方式承担征发、保护、迁移或宗教活动；现存来源只能照亮其中一部分，不能替沉默者制造统一的经验。

\eventanchor{JH-LEDGER-006-2}{元康六年至九年（296年至299年）·“氐族齐万年在关中起兵，西晋多次征讨并长期调动关陇军粮，至299年叛乱被平定”的命令、联盟或地方响应}账本记录西晋与关陇氐羌集团在元康六年至九年（296年至299年）发生““氐族齐万年在关中起兵，西晋多次征讨并长期调动关陇军粮，至299年叛乱被平定”的命令、联盟或地方响应”。条目把背景概括为当时的权力、资源与信息约束推动相关过程展开，并指出这一过程为“氐族齐万年在关中起兵，西晋多次征讨并长期调动关陇军粮，至299年叛乱被平定”提供可观察的条件。这个节点进入区域社会时，要经由文书、道路、军队、市场、寺院和家庭等不同中介；因此，政治行动的宣布、地方接收和日常生活的改变不能被视作同一时刻完成。

本条依据《晋书·惠帝纪》；《晋书·张方传》；《晋书·江统传》；《资治通鉴·晋纪》，证据提示为“核心事项已有既有账本来源；本分解节点的参与者、执行范围和时间先后仍须逐案核验”。这些材料可支持年序、官方决定、战事或特定地点的线索，却难以直接统计所有人的财富、迁徙、信仰和动机。更稳妥的读法是区分诏令与实施、军报与人口、行政称谓与自我认同，并让地方材料的缺环限制解释的范围。

由此可见，该事件不是孤立的年表点，而是资源、信息与权力关系重新配置的一环。不同地区的官员、社群和家庭可能以不同方式承担征发、保护、迁移或宗教活动；现存来源只能照亮其中一部分，不能替沉默者制造统一的经验。

\eventanchor{JH-LEDGER-006-3}{元康六年至九年（296年至299年）·氐族齐万年在关中起兵，西晋多次征讨并长期调动关陇军粮，至299年叛乱被平定}账本记录西晋与关陇氐羌集团在元康六年至九年（296年至299年）发生“氐族齐万年在关中起兵，西晋多次征讨并长期调动关陇军粮，至299年叛乱被平定”。条目把背景概括为“氐族齐万年在关中起兵，西晋多次征讨并长期调动关陇军粮，至299年叛乱被平定”发生的政治与区域条件共同作用，并指出该节点改变后续资源、联盟或制度处置的条件。这个节点进入区域社会时，要经由文书、道路、军队、市场、寺院和家庭等不同中介；因此，政治行动的宣布、地方接收和日常生活的改变不能被视作同一时刻完成。

本条依据《晋书·惠帝纪》；《晋书·张方传》；《晋书·江统传》；《资治通鉴·晋纪》，证据提示为“核心事项已有既有账本来源；本分解节点的参与者、执行范围和时间先后仍须逐案核验”。这些材料可支持年序、官方决定、战事或特定地点的线索，却难以直接统计所有人的财富、迁徙、信仰和动机。更稳妥的读法是区分诏令与实施、军报与人口、行政称谓与自我认同，并让地方材料的缺环限制解释的范围。

由此可见，该事件不是孤立的年表点，而是资源、信息与权力关系重新配置的一环。不同地区的官员、社群和家庭可能以不同方式承担征发、保护、迁移或宗教活动；现存来源只能照亮其中一部分，不能替沉默者制造统一的经验。

\eventanchor{JH-LEDGER-006-4}{元康六年至九年（296年至299年）·“氐族齐万年在关中起兵，西晋多次征讨并长期调动关陇军粮，至299年叛乱被平定”的执行中的空间与人员处置}账本记录西晋与关陇氐羌集团在元康六年至九年（296年至299年）发生““氐族齐万年在关中起兵，西晋多次征讨并长期调动关陇军粮，至299年叛乱被平定”的执行中的空间与人员处置”。条目把背景概括为“氐族齐万年在关中起兵，西晋多次征讨并长期调动关陇军粮，至299年叛乱被平定”发生的政治与区域条件共同作用，并指出该节点改变后续资源、联盟或制度处置的条件。这个节点进入区域社会时，要经由文书、道路、军队、市场、寺院和家庭等不同中介；因此，政治行动的宣布、地方接收和日常生活的改变不能被视作同一时刻完成。

本条依据《晋书·惠帝纪》；《晋书·张方传》；《晋书·江统传》；《资治通鉴·晋纪》，证据提示为“核心事项已有既有账本来源；本分解节点的参与者、执行范围和时间先后仍须逐案核验”。这些材料可支持年序、官方决定、战事或特定地点的线索，却难以直接统计所有人的财富、迁徙、信仰和动机。更稳妥的读法是区分诏令与实施、军报与人口、行政称谓与自我认同，并让地方材料的缺环限制解释的范围。

由此可见，该事件不是孤立的年表点，而是资源、信息与权力关系重新配置的一环。不同地区的官员、社群和家庭可能以不同方式承担征发、保护、迁移或宗教活动；现存来源只能照亮其中一部分，不能替沉默者制造统一的经验。

\eventanchor{JH-LEDGER-006-5}{元康六年至九年（296年至299年）·“氐族齐万年在关中起兵，西晋多次征讨并长期调动关陇军粮，至299年叛乱被平定”的后续制度或军政调整}账本记录西晋与关陇氐羌集团在元康六年至九年（296年至299年）发生““氐族齐万年在关中起兵，西晋多次征讨并长期调动关陇军粮，至299年叛乱被平定”的后续制度或军政调整”。条目把背景概括为核心节点发生后仍须处理其遗留问题，并指出后续影响须在不同地区和材料类型中分别核验。这个节点进入区域社会时，要经由文书、道路、军队、市场、寺院和家庭等不同中介；因此，政治行动的宣布、地方接收和日常生活的改变不能被视作同一时刻完成。

本条依据《晋书·惠帝纪》；《晋书·张方传》；《晋书·江统传》；《资治通鉴·晋纪》，证据提示为“核心事项已有既有账本来源；本分解节点的参与者、执行范围和时间先后仍须逐案核验”。这些材料可支持年序、官方决定、战事或特定地点的线索，却难以直接统计所有人的财富、迁徙、信仰和动机。更稳妥的读法是区分诏令与实施、军报与人口、行政称谓与自我认同，并让地方材料的缺环限制解释的范围。

由此可见，该事件不是孤立的年表点，而是资源、信息与权力关系重新配置的一环。不同地区的官员、社群和家庭可能以不同方式承担征发、保护、迁移或宗教活动；现存来源只能照亮其中一部分，不能替沉默者制造统一的经验。

\eventanchor{JH-LEDGER-006-6}{元康六年至九年（296年至299年）·“氐族齐万年在关中起兵，西晋多次征讨并长期调动关陇军粮，至299年叛乱被平定”的异源材料与影响范围的复核}账本记录西晋与关陇氐羌集团在元康六年至九年（296年至299年）发生““氐族齐万年在关中起兵，西晋多次征讨并长期调动关陇军粮，至299年叛乱被平定”的异源材料与影响范围的复核”。条目把背景概括为核心节点发生后仍须处理其遗留问题，并指出后续影响须在不同地区和材料类型中分别核验。这个节点进入区域社会时，要经由文书、道路、军队、市场、寺院和家庭等不同中介；因此，政治行动的宣布、地方接收和日常生活的改变不能被视作同一时刻完成。

本条依据《晋书·惠帝纪》；《晋书·张方传》；《晋书·江统传》；《资治通鉴·晋纪》，证据提示为“核心事项已有既有账本来源；本分解节点的参与者、执行范围和时间先后仍须逐案核验”。这些材料可支持年序、官方决定、战事或特定地点的线索，却难以直接统计所有人的财富、迁徙、信仰和动机。更稳妥的读法是区分诏令与实施、军报与人口、行政称谓与自我认同，并让地方材料的缺环限制解释的范围。

由此可见，该事件不是孤立的年表点，而是资源、信息与权力关系重新配置的一环。不同地区的官员、社群和家庭可能以不同方式承担征发、保护、迁移或宗教活动；现存来源只能照亮其中一部分，不能替沉默者制造统一的经验。

\eventanchor{JH-LEDGER-007-1}{永康元年（300年）·“贾后以谋反罪废杀太子司马遹，皇位继承与宫廷联盟由此公开破裂”的前期动员与资源准备}账本记录西晋在永康元年（300年）发生““贾后以谋反罪废杀太子司马遹，皇位继承与宫廷联盟由此公开破裂”的前期动员与资源准备”。条目把背景概括为当时的权力、资源与信息约束推动相关过程展开，并指出这一过程为“贾后以谋反罪废杀太子司马遹，皇位继承与宫廷联盟由此公开破裂”提供可观察的条件。这个节点进入区域社会时，要经由文书、道路、军队、市场、寺院和家庭等不同中介；因此，政治行动的宣布、地方接收和日常生活的改变不能被视作同一时刻完成。

本条依据《晋书·惠帝纪》；《晋书·愍怀太子遹传》；《晋书·贾后传》，证据提示为“核心事项已有既有账本来源；本分解节点的参与者、执行范围和时间先后仍须逐案核验”。这些材料可支持年序、官方决定、战事或特定地点的线索，却难以直接统计所有人的财富、迁徙、信仰和动机。更稳妥的读法是区分诏令与实施、军报与人口、行政称谓与自我认同，并让地方材料的缺环限制解释的范围。

由此可见，该事件不是孤立的年表点，而是资源、信息与权力关系重新配置的一环。不同地区的官员、社群和家庭可能以不同方式承担征发、保护、迁移或宗教活动；现存来源只能照亮其中一部分，不能替沉默者制造统一的经验。

\eventanchor{JH-LEDGER-007-2}{永康元年（300年）·“贾后以谋反罪废杀太子司马遹，皇位继承与宫廷联盟由此公开破裂”的命令、联盟或地方响应}账本记录西晋在永康元年（300年）发生““贾后以谋反罪废杀太子司马遹，皇位继承与宫廷联盟由此公开破裂”的命令、联盟或地方响应”。条目把背景概括为当时的权力、资源与信息约束推动相关过程展开，并指出这一过程为“贾后以谋反罪废杀太子司马遹，皇位继承与宫廷联盟由此公开破裂”提供可观察的条件。这个节点进入区域社会时，要经由文书、道路、军队、市场、寺院和家庭等不同中介；因此，政治行动的宣布、地方接收和日常生活的改变不能被视作同一时刻完成。

本条依据《晋书·惠帝纪》；《晋书·愍怀太子遹传》；《晋书·贾后传》，证据提示为“核心事项已有既有账本来源；本分解节点的参与者、执行范围和时间先后仍须逐案核验”。这些材料可支持年序、官方决定、战事或特定地点的线索，却难以直接统计所有人的财富、迁徙、信仰和动机。更稳妥的读法是区分诏令与实施、军报与人口、行政称谓与自我认同，并让地方材料的缺环限制解释的范围。

由此可见，该事件不是孤立的年表点，而是资源、信息与权力关系重新配置的一环。不同地区的官员、社群和家庭可能以不同方式承担征发、保护、迁移或宗教活动；现存来源只能照亮其中一部分，不能替沉默者制造统一的经验。

\eventanchor{JH-LEDGER-007-3}{永康元年（300年）·贾后以谋反罪废杀太子司马遹，皇位继承与宫廷联盟由此公开破裂}账本记录西晋在永康元年（300年）发生“贾后以谋反罪废杀太子司马遹，皇位继承与宫廷联盟由此公开破裂”。条目把背景概括为“贾后以谋反罪废杀太子司马遹，皇位继承与宫廷联盟由此公开破裂”发生的政治与区域条件共同作用，并指出该节点改变后续资源、联盟或制度处置的条件。这个节点进入区域社会时，要经由文书、道路、军队、市场、寺院和家庭等不同中介；因此，政治行动的宣布、地方接收和日常生活的改变不能被视作同一时刻完成。

本条依据《晋书·惠帝纪》；《晋书·愍怀太子遹传》；《晋书·贾后传》，证据提示为“核心事项已有既有账本来源；本分解节点的参与者、执行范围和时间先后仍须逐案核验”。这些材料可支持年序、官方决定、战事或特定地点的线索，却难以直接统计所有人的财富、迁徙、信仰和动机。更稳妥的读法是区分诏令与实施、军报与人口、行政称谓与自我认同，并让地方材料的缺环限制解释的范围。

由此可见，该事件不是孤立的年表点，而是资源、信息与权力关系重新配置的一环。不同地区的官员、社群和家庭可能以不同方式承担征发、保护、迁移或宗教活动；现存来源只能照亮其中一部分，不能替沉默者制造统一的经验。

\eventanchor{JH-LEDGER-007-4}{永康元年（300年）·“贾后以谋反罪废杀太子司马遹，皇位继承与宫廷联盟由此公开破裂”的执行中的空间与人员处置}账本记录西晋在永康元年（300年）发生““贾后以谋反罪废杀太子司马遹，皇位继承与宫廷联盟由此公开破裂”的执行中的空间与人员处置”。条目把背景概括为“贾后以谋反罪废杀太子司马遹，皇位继承与宫廷联盟由此公开破裂”发生的政治与区域条件共同作用，并指出该节点改变后续资源、联盟或制度处置的条件。这个节点进入区域社会时，要经由文书、道路、军队、市场、寺院和家庭等不同中介；因此，政治行动的宣布、地方接收和日常生活的改变不能被视作同一时刻完成。

本条依据《晋书·惠帝纪》；《晋书·愍怀太子遹传》；《晋书·贾后传》，证据提示为“核心事项已有既有账本来源；本分解节点的参与者、执行范围和时间先后仍须逐案核验”。这些材料可支持年序、官方决定、战事或特定地点的线索，却难以直接统计所有人的财富、迁徙、信仰和动机。更稳妥的读法是区分诏令与实施、军报与人口、行政称谓与自我认同，并让地方材料的缺环限制解释的范围。

由此可见，该事件不是孤立的年表点，而是资源、信息与权力关系重新配置的一环。不同地区的官员、社群和家庭可能以不同方式承担征发、保护、迁移或宗教活动；现存来源只能照亮其中一部分，不能替沉默者制造统一的经验。

\eventanchor{JH-LEDGER-007-5}{永康元年（300年）·“贾后以谋反罪废杀太子司马遹，皇位继承与宫廷联盟由此公开破裂”的后续制度或军政调整}账本记录西晋在永康元年（300年）发生““贾后以谋反罪废杀太子司马遹，皇位继承与宫廷联盟由此公开破裂”的后续制度或军政调整”。条目把背景概括为核心节点发生后仍须处理其遗留问题，并指出后续影响须在不同地区和材料类型中分别核验。这个节点进入区域社会时，要经由文书、道路、军队、市场、寺院和家庭等不同中介；因此，政治行动的宣布、地方接收和日常生活的改变不能被视作同一时刻完成。

本条依据《晋书·惠帝纪》；《晋书·愍怀太子遹传》；《晋书·贾后传》，证据提示为“核心事项已有既有账本来源；本分解节点的参与者、执行范围和时间先后仍须逐案核验”。这些材料可支持年序、官方决定、战事或特定地点的线索，却难以直接统计所有人的财富、迁徙、信仰和动机。更稳妥的读法是区分诏令与实施、军报与人口、行政称谓与自我认同，并让地方材料的缺环限制解释的范围。

由此可见，该事件不是孤立的年表点，而是资源、信息与权力关系重新配置的一环。不同地区的官员、社群和家庭可能以不同方式承担征发、保护、迁移或宗教活动；现存来源只能照亮其中一部分，不能替沉默者制造统一的经验。

\eventanchor{JH-LEDGER-007-6}{永康元年（300年）·“贾后以谋反罪废杀太子司马遹，皇位继承与宫廷联盟由此公开破裂”的异源材料与影响范围的复核}账本记录西晋在永康元年（300年）发生““贾后以谋反罪废杀太子司马遹，皇位继承与宫廷联盟由此公开破裂”的异源材料与影响范围的复核”。条目把背景概括为核心节点发生后仍须处理其遗留问题，并指出后续影响须在不同地区和材料类型中分别核验。这个节点进入区域社会时，要经由文书、道路、军队、市场、寺院和家庭等不同中介；因此，政治行动的宣布、地方接收和日常生活的改变不能被视作同一时刻完成。

本条依据《晋书·惠帝纪》；《晋书·愍怀太子遹传》；《晋书·贾后传》，证据提示为“核心事项已有既有账本来源；本分解节点的参与者、执行范围和时间先后仍须逐案核验”。这些材料可支持年序、官方决定、战事或特定地点的线索，却难以直接统计所有人的财富、迁徙、信仰和动机。更稳妥的读法是区分诏令与实施、军报与人口、行政称谓与自我认同，并让地方材料的缺环限制解释的范围。

由此可见，该事件不是孤立的年表点，而是资源、信息与权力关系重新配置的一环。不同地区的官员、社群和家庭可能以不同方式承担征发、保护、迁移或宗教活动；现存来源只能照亮其中一部分，不能替沉默者制造统一的经验。

\eventanchor{JH-LEDGER-008-1}{永宁元年（301年）·“赵王司马伦杀贾后并废惠帝自立为帝，短暂改变西晋帝统”的前期动员与资源准备}账本记录西晋在永宁元年（301年）发生““赵王司马伦杀贾后并废惠帝自立为帝，短暂改变西晋帝统”的前期动员与资源准备”。条目把背景概括为当时的权力、资源与信息约束推动相关过程展开，并指出这一过程为“赵王司马伦杀贾后并废惠帝自立为帝，短暂改变西晋帝统”提供可观察的条件。这个节点进入区域社会时，要经由文书、道路、军队、市场、寺院和家庭等不同中介；因此，政治行动的宣布、地方接收和日常生活的改变不能被视作同一时刻完成。

本条依据《晋书·惠帝纪》；《晋书·赵王伦传》；《晋书·贾后传》，证据提示为“核心事项已有既有账本来源；本分解节点的参与者、执行范围和时间先后仍须逐案核验”。这些材料可支持年序、官方决定、战事或特定地点的线索，却难以直接统计所有人的财富、迁徙、信仰和动机。更稳妥的读法是区分诏令与实施、军报与人口、行政称谓与自我认同，并让地方材料的缺环限制解释的范围。

由此可见，该事件不是孤立的年表点，而是资源、信息与权力关系重新配置的一环。不同地区的官员、社群和家庭可能以不同方式承担征发、保护、迁移或宗教活动；现存来源只能照亮其中一部分，不能替沉默者制造统一的经验。

\eventanchor{JH-LEDGER-008-2}{永宁元年（301年）·“赵王司马伦杀贾后并废惠帝自立为帝，短暂改变西晋帝统”的命令、联盟或地方响应}账本记录西晋在永宁元年（301年）发生““赵王司马伦杀贾后并废惠帝自立为帝，短暂改变西晋帝统”的命令、联盟或地方响应”。条目把背景概括为当时的权力、资源与信息约束推动相关过程展开，并指出这一过程为“赵王司马伦杀贾后并废惠帝自立为帝，短暂改变西晋帝统”提供可观察的条件。这个节点进入区域社会时，要经由文书、道路、军队、市场、寺院和家庭等不同中介；因此，政治行动的宣布、地方接收和日常生活的改变不能被视作同一时刻完成。

本条依据《晋书·惠帝纪》；《晋书·赵王伦传》；《晋书·贾后传》，证据提示为“核心事项已有既有账本来源；本分解节点的参与者、执行范围和时间先后仍须逐案核验”。这些材料可支持年序、官方决定、战事或特定地点的线索，却难以直接统计所有人的财富、迁徙、信仰和动机。更稳妥的读法是区分诏令与实施、军报与人口、行政称谓与自我认同，并让地方材料的缺环限制解释的范围。

由此可见，该事件不是孤立的年表点，而是资源、信息与权力关系重新配置的一环。不同地区的官员、社群和家庭可能以不同方式承担征发、保护、迁移或宗教活动；现存来源只能照亮其中一部分，不能替沉默者制造统一的经验。

\eventanchor{JH-LEDGER-008-3}{永宁元年（301年）·赵王司马伦杀贾后并废惠帝自立为帝，短暂改变西晋帝统}账本记录西晋在永宁元年（301年）发生“赵王司马伦杀贾后并废惠帝自立为帝，短暂改变西晋帝统”。条目把背景概括为“赵王司马伦杀贾后并废惠帝自立为帝，短暂改变西晋帝统”发生的政治与区域条件共同作用，并指出该节点改变后续资源、联盟或制度处置的条件。这个节点进入区域社会时，要经由文书、道路、军队、市场、寺院和家庭等不同中介；因此，政治行动的宣布、地方接收和日常生活的改变不能被视作同一时刻完成。

本条依据《晋书·惠帝纪》；《晋书·赵王伦传》；《晋书·贾后传》，证据提示为“核心事项已有既有账本来源；本分解节点的参与者、执行范围和时间先后仍须逐案核验”。这些材料可支持年序、官方决定、战事或特定地点的线索，却难以直接统计所有人的财富、迁徙、信仰和动机。更稳妥的读法是区分诏令与实施、军报与人口、行政称谓与自我认同，并让地方材料的缺环限制解释的范围。

由此可见，该事件不是孤立的年表点，而是资源、信息与权力关系重新配置的一环。不同地区的官员、社群和家庭可能以不同方式承担征发、保护、迁移或宗教活动；现存来源只能照亮其中一部分，不能替沉默者制造统一的经验。

\eventanchor{JH-LEDGER-008-4}{永宁元年（301年）·“赵王司马伦杀贾后并废惠帝自立为帝，短暂改变西晋帝统”的执行中的空间与人员处置}账本记录西晋在永宁元年（301年）发生““赵王司马伦杀贾后并废惠帝自立为帝，短暂改变西晋帝统”的执行中的空间与人员处置”。条目把背景概括为“赵王司马伦杀贾后并废惠帝自立为帝，短暂改变西晋帝统”发生的政治与区域条件共同作用，并指出该节点改变后续资源、联盟或制度处置的条件。这个节点进入区域社会时，要经由文书、道路、军队、市场、寺院和家庭等不同中介；因此，政治行动的宣布、地方接收和日常生活的改变不能被视作同一时刻完成。

本条依据《晋书·惠帝纪》；《晋书·赵王伦传》；《晋书·贾后传》，证据提示为“核心事项已有既有账本来源；本分解节点的参与者、执行范围和时间先后仍须逐案核验”。这些材料可支持年序、官方决定、战事或特定地点的线索，却难以直接统计所有人的财富、迁徙、信仰和动机。更稳妥的读法是区分诏令与实施、军报与人口、行政称谓与自我认同，并让地方材料的缺环限制解释的范围。

由此可见，该事件不是孤立的年表点，而是资源、信息与权力关系重新配置的一环。不同地区的官员、社群和家庭可能以不同方式承担征发、保护、迁移或宗教活动；现存来源只能照亮其中一部分，不能替沉默者制造统一的经验。

\subsection{边疆、战争与移动}

\eventanchor{JH-LEDGER-008-5}{永宁元年（301年）·“赵王司马伦杀贾后并废惠帝自立为帝，短暂改变西晋帝统”的后续制度或军政调整}账本记录西晋在永宁元年（301年）发生““赵王司马伦杀贾后并废惠帝自立为帝，短暂改变西晋帝统”的后续制度或军政调整”。条目把背景概括为核心节点发生后仍须处理其遗留问题，并指出后续影响须在不同地区和材料类型中分别核验。这个节点进入区域社会时，要经由文书、道路、军队、市场、寺院和家庭等不同中介；因此，政治行动的宣布、地方接收和日常生活的改变不能被视作同一时刻完成。

本条依据《晋书·惠帝纪》；《晋书·赵王伦传》；《晋书·贾后传》，证据提示为“核心事项已有既有账本来源；本分解节点的参与者、执行范围和时间先后仍须逐案核验”。这些材料可支持年序、官方决定、战事或特定地点的线索，却难以直接统计所有人的财富、迁徙、信仰和动机。更稳妥的读法是区分诏令与实施、军报与人口、行政称谓与自我认同，并让地方材料的缺环限制解释的范围。

由此可见，该事件不是孤立的年表点，而是资源、信息与权力关系重新配置的一环。不同地区的官员、社群和家庭可能以不同方式承担征发、保护、迁移或宗教活动；现存来源只能照亮其中一部分，不能替沉默者制造统一的经验。

\eventanchor{JH-LEDGER-008-6}{永宁元年（301年）·“赵王司马伦杀贾后并废惠帝自立为帝，短暂改变西晋帝统”的异源材料与影响范围的复核}账本记录西晋在永宁元年（301年）发生““赵王司马伦杀贾后并废惠帝自立为帝，短暂改变西晋帝统”的异源材料与影响范围的复核”。条目把背景概括为核心节点发生后仍须处理其遗留问题，并指出后续影响须在不同地区和材料类型中分别核验。这个节点进入区域社会时，要经由文书、道路、军队、市场、寺院和家庭等不同中介；因此，政治行动的宣布、地方接收和日常生活的改变不能被视作同一时刻完成。

本条依据《晋书·惠帝纪》；《晋书·赵王伦传》；《晋书·贾后传》，证据提示为“核心事项已有既有账本来源；本分解节点的参与者、执行范围和时间先后仍须逐案核验”。这些材料可支持年序、官方决定、战事或特定地点的线索，却难以直接统计所有人的财富、迁徙、信仰和动机。更稳妥的读法是区分诏令与实施、军报与人口、行政称谓与自我认同，并让地方材料的缺环限制解释的范围。

由此可见，该事件不是孤立的年表点，而是资源、信息与权力关系重新配置的一环。不同地区的官员、社群和家庭可能以不同方式承担征发、保护、迁移或宗教活动；现存来源只能照亮其中一部分，不能替沉默者制造统一的经验。

\eventanchor{JH-LEDGER-009-1}{永宁元年（301年）·“齐王司马冏等起兵推翻司马伦，复立惠帝并以大司马辅政”的前期动员与资源准备}账本记录西晋在永宁元年（301年）发生““齐王司马冏等起兵推翻司马伦，复立惠帝并以大司马辅政”的前期动员与资源准备”。条目把背景概括为当时的权力、资源与信息约束推动相关过程展开，并指出这一过程为“齐王司马冏等起兵推翻司马伦，复立惠帝并以大司马辅政”提供可观察的条件。这个节点进入区域社会时，要经由文书、道路、军队、市场、寺院和家庭等不同中介；因此，政治行动的宣布、地方接收和日常生活的改变不能被视作同一时刻完成。

本条依据《晋书·惠帝纪》；《晋书·齐王冏传》；《晋书·赵王伦传》；《资治通鉴·晋纪》，证据提示为“核心事项已有既有账本来源；本分解节点的参与者、执行范围和时间先后仍须逐案核验”。这些材料可支持年序、官方决定、战事或特定地点的线索，却难以直接统计所有人的财富、迁徙、信仰和动机。更稳妥的读法是区分诏令与实施、军报与人口、行政称谓与自我认同，并让地方材料的缺环限制解释的范围。

由此可见，该事件不是孤立的年表点，而是资源、信息与权力关系重新配置的一环。不同地区的官员、社群和家庭可能以不同方式承担征发、保护、迁移或宗教活动；现存来源只能照亮其中一部分，不能替沉默者制造统一的经验。

\eventanchor{JH-LEDGER-009-2}{永宁元年（301年）·“齐王司马冏等起兵推翻司马伦，复立惠帝并以大司马辅政”的命令、联盟或地方响应}账本记录西晋在永宁元年（301年）发生““齐王司马冏等起兵推翻司马伦，复立惠帝并以大司马辅政”的命令、联盟或地方响应”。条目把背景概括为当时的权力、资源与信息约束推动相关过程展开，并指出这一过程为“齐王司马冏等起兵推翻司马伦，复立惠帝并以大司马辅政”提供可观察的条件。这个节点进入区域社会时，要经由文书、道路、军队、市场、寺院和家庭等不同中介；因此，政治行动的宣布、地方接收和日常生活的改变不能被视作同一时刻完成。

本条依据《晋书·惠帝纪》；《晋书·齐王冏传》；《晋书·赵王伦传》；《资治通鉴·晋纪》，证据提示为“核心事项已有既有账本来源；本分解节点的参与者、执行范围和时间先后仍须逐案核验”。这些材料可支持年序、官方决定、战事或特定地点的线索，却难以直接统计所有人的财富、迁徙、信仰和动机。更稳妥的读法是区分诏令与实施、军报与人口、行政称谓与自我认同，并让地方材料的缺环限制解释的范围。

由此可见，该事件不是孤立的年表点，而是资源、信息与权力关系重新配置的一环。不同地区的官员、社群和家庭可能以不同方式承担征发、保护、迁移或宗教活动；现存来源只能照亮其中一部分，不能替沉默者制造统一的经验。

\eventanchor{JH-LEDGER-009-3}{永宁元年（301年）·齐王司马冏等起兵推翻司马伦，复立惠帝并以大司马辅政}账本记录西晋在永宁元年（301年）发生“齐王司马冏等起兵推翻司马伦，复立惠帝并以大司马辅政”。条目把背景概括为“齐王司马冏等起兵推翻司马伦，复立惠帝并以大司马辅政”发生的政治与区域条件共同作用，并指出该节点改变后续资源、联盟或制度处置的条件。这个节点进入区域社会时，要经由文书、道路、军队、市场、寺院和家庭等不同中介；因此，政治行动的宣布、地方接收和日常生活的改变不能被视作同一时刻完成。

本条依据《晋书·惠帝纪》；《晋书·齐王冏传》；《晋书·赵王伦传》；《资治通鉴·晋纪》，证据提示为“核心事项已有既有账本来源；本分解节点的参与者、执行范围和时间先后仍须逐案核验”。这些材料可支持年序、官方决定、战事或特定地点的线索，却难以直接统计所有人的财富、迁徙、信仰和动机。更稳妥的读法是区分诏令与实施、军报与人口、行政称谓与自我认同，并让地方材料的缺环限制解释的范围。

由此可见，该事件不是孤立的年表点，而是资源、信息与权力关系重新配置的一环。不同地区的官员、社群和家庭可能以不同方式承担征发、保护、迁移或宗教活动；现存来源只能照亮其中一部分，不能替沉默者制造统一的经验。

\eventanchor{JH-LEDGER-009-4}{永宁元年（301年）·“齐王司马冏等起兵推翻司马伦，复立惠帝并以大司马辅政”的执行中的空间与人员处置}账本记录西晋在永宁元年（301年）发生““齐王司马冏等起兵推翻司马伦，复立惠帝并以大司马辅政”的执行中的空间与人员处置”。条目把背景概括为“齐王司马冏等起兵推翻司马伦，复立惠帝并以大司马辅政”发生的政治与区域条件共同作用，并指出该节点改变后续资源、联盟或制度处置的条件。这个节点进入区域社会时，要经由文书、道路、军队、市场、寺院和家庭等不同中介；因此，政治行动的宣布、地方接收和日常生活的改变不能被视作同一时刻完成。

本条依据《晋书·惠帝纪》；《晋书·齐王冏传》；《晋书·赵王伦传》；《资治通鉴·晋纪》，证据提示为“核心事项已有既有账本来源；本分解节点的参与者、执行范围和时间先后仍须逐案核验”。这些材料可支持年序、官方决定、战事或特定地点的线索，却难以直接统计所有人的财富、迁徙、信仰和动机。更稳妥的读法是区分诏令与实施、军报与人口、行政称谓与自我认同，并让地方材料的缺环限制解释的范围。

由此可见，该事件不是孤立的年表点，而是资源、信息与权力关系重新配置的一环。不同地区的官员、社群和家庭可能以不同方式承担征发、保护、迁移或宗教活动；现存来源只能照亮其中一部分，不能替沉默者制造统一的经验。

\eventanchor{JH-LEDGER-009-5}{永宁元年（301年）·“齐王司马冏等起兵推翻司马伦，复立惠帝并以大司马辅政”的后续制度或军政调整}账本记录西晋在永宁元年（301年）发生““齐王司马冏等起兵推翻司马伦，复立惠帝并以大司马辅政”的后续制度或军政调整”。条目把背景概括为核心节点发生后仍须处理其遗留问题，并指出后续影响须在不同地区和材料类型中分别核验。这个节点进入区域社会时，要经由文书、道路、军队、市场、寺院和家庭等不同中介；因此，政治行动的宣布、地方接收和日常生活的改变不能被视作同一时刻完成。

本条依据《晋书·惠帝纪》；《晋书·齐王冏传》；《晋书·赵王伦传》；《资治通鉴·晋纪》，证据提示为“核心事项已有既有账本来源；本分解节点的参与者、执行范围和时间先后仍须逐案核验”。这些材料可支持年序、官方决定、战事或特定地点的线索，却难以直接统计所有人的财富、迁徙、信仰和动机。更稳妥的读法是区分诏令与实施、军报与人口、行政称谓与自我认同，并让地方材料的缺环限制解释的范围。

由此可见，该事件不是孤立的年表点，而是资源、信息与权力关系重新配置的一环。不同地区的官员、社群和家庭可能以不同方式承担征发、保护、迁移或宗教活动；现存来源只能照亮其中一部分，不能替沉默者制造统一的经验。

\eventanchor{JH-LEDGER-009-6}{永宁元年（301年）·“齐王司马冏等起兵推翻司马伦，复立惠帝并以大司马辅政”的异源材料与影响范围的复核}账本记录西晋在永宁元年（301年）发生““齐王司马冏等起兵推翻司马伦，复立惠帝并以大司马辅政”的异源材料与影响范围的复核”。条目把背景概括为核心节点发生后仍须处理其遗留问题，并指出后续影响须在不同地区和材料类型中分别核验。这个节点进入区域社会时，要经由文书、道路、军队、市场、寺院和家庭等不同中介；因此，政治行动的宣布、地方接收和日常生活的改变不能被视作同一时刻完成。

本条依据《晋书·惠帝纪》；《晋书·齐王冏传》；《晋书·赵王伦传》；《资治通鉴·晋纪》，证据提示为“核心事项已有既有账本来源；本分解节点的参与者、执行范围和时间先后仍须逐案核验”。这些材料可支持年序、官方决定、战事或特定地点的线索，却难以直接统计所有人的财富、迁徙、信仰和动机。更稳妥的读法是区分诏令与实施、军报与人口、行政称谓与自我认同，并让地方材料的缺环限制解释的范围。

由此可见，该事件不是孤立的年表点，而是资源、信息与权力关系重新配置的一环。不同地区的官员、社群和家庭可能以不同方式承担征发、保护、迁移或宗教活动；现存来源只能照亮其中一部分，不能替沉默者制造统一的经验。

\eventanchor{JH-LEDGER-010-1}{太安元年（302年）·“河间王司马颙与成都王司马颖起兵，长沙王司马乂在洛阳杀齐王司马冏并接掌朝政”的前期动员与资源准备}账本记录西晋在太安元年（302年）发生““河间王司马颙与成都王司马颖起兵，长沙王司马乂在洛阳杀齐王司马冏并接掌朝政”的前期动员与资源准备”。条目把背景概括为当时的权力、资源与信息约束推动相关过程展开，并指出这一过程为“河间王司马颙与成都王司马颖起兵，长沙王司马乂在洛阳杀齐王司马冏并接掌朝政”提供可观察的条件。这个节点进入区域社会时，要经由文书、道路、军队、市场、寺院和家庭等不同中介；因此，政治行动的宣布、地方接收和日常生活的改变不能被视作同一时刻完成。

本条依据《晋书·惠帝纪》；《晋书·齐王冏传》；《晋书·长沙厉王乂传》；《晋书·成都王颖传》，证据提示为“核心事项已有既有账本来源；本分解节点的参与者、执行范围和时间先后仍须逐案核验”。这些材料可支持年序、官方决定、战事或特定地点的线索，却难以直接统计所有人的财富、迁徙、信仰和动机。更稳妥的读法是区分诏令与实施、军报与人口、行政称谓与自我认同，并让地方材料的缺环限制解释的范围。

由此可见，该事件不是孤立的年表点，而是资源、信息与权力关系重新配置的一环。不同地区的官员、社群和家庭可能以不同方式承担征发、保护、迁移或宗教活动；现存来源只能照亮其中一部分，不能替沉默者制造统一的经验。

\eventanchor{JH-LEDGER-010-2}{太安元年（302年）·“河间王司马颙与成都王司马颖起兵，长沙王司马乂在洛阳杀齐王司马冏并接掌朝政”的命令、联盟或地方响应}账本记录西晋在太安元年（302年）发生““河间王司马颙与成都王司马颖起兵，长沙王司马乂在洛阳杀齐王司马冏并接掌朝政”的命令、联盟或地方响应”。条目把背景概括为当时的权力、资源与信息约束推动相关过程展开，并指出这一过程为“河间王司马颙与成都王司马颖起兵，长沙王司马乂在洛阳杀齐王司马冏并接掌朝政”提供可观察的条件。这个节点进入区域社会时，要经由文书、道路、军队、市场、寺院和家庭等不同中介；因此，政治行动的宣布、地方接收和日常生活的改变不能被视作同一时刻完成。

本条依据《晋书·惠帝纪》；《晋书·齐王冏传》；《晋书·长沙厉王乂传》；《晋书·成都王颖传》，证据提示为“核心事项已有既有账本来源；本分解节点的参与者、执行范围和时间先后仍须逐案核验”。这些材料可支持年序、官方决定、战事或特定地点的线索，却难以直接统计所有人的财富、迁徙、信仰和动机。更稳妥的读法是区分诏令与实施、军报与人口、行政称谓与自我认同，并让地方材料的缺环限制解释的范围。

由此可见，该事件不是孤立的年表点，而是资源、信息与权力关系重新配置的一环。不同地区的官员、社群和家庭可能以不同方式承担征发、保护、迁移或宗教活动；现存来源只能照亮其中一部分，不能替沉默者制造统一的经验。

\eventanchor{JH-LEDGER-010-3}{太安元年（302年）·河间王司马颙与成都王司马颖起兵，长沙王司马乂在洛阳杀齐王司马冏并接掌朝政}账本记录西晋在太安元年（302年）发生“河间王司马颙与成都王司马颖起兵，长沙王司马乂在洛阳杀齐王司马冏并接掌朝政”。条目把背景概括为“河间王司马颙与成都王司马颖起兵，长沙王司马乂在洛阳杀齐王司马冏并接掌朝政”发生的政治与区域条件共同作用，并指出该节点改变后续资源、联盟或制度处置的条件。这个节点进入区域社会时，要经由文书、道路、军队、市场、寺院和家庭等不同中介；因此，政治行动的宣布、地方接收和日常生活的改变不能被视作同一时刻完成。

本条依据《晋书·惠帝纪》；《晋书·齐王冏传》；《晋书·长沙厉王乂传》；《晋书·成都王颖传》，证据提示为“核心事项已有既有账本来源；本分解节点的参与者、执行范围和时间先后仍须逐案核验”。这些材料可支持年序、官方决定、战事或特定地点的线索，却难以直接统计所有人的财富、迁徙、信仰和动机。更稳妥的读法是区分诏令与实施、军报与人口、行政称谓与自我认同，并让地方材料的缺环限制解释的范围。

由此可见，该事件不是孤立的年表点，而是资源、信息与权力关系重新配置的一环。不同地区的官员、社群和家庭可能以不同方式承担征发、保护、迁移或宗教活动；现存来源只能照亮其中一部分，不能替沉默者制造统一的经验。

\eventanchor{JH-LEDGER-010-4}{太安元年（302年）·“河间王司马颙与成都王司马颖起兵，长沙王司马乂在洛阳杀齐王司马冏并接掌朝政”的执行中的空间与人员处置}账本记录西晋在太安元年（302年）发生““河间王司马颙与成都王司马颖起兵，长沙王司马乂在洛阳杀齐王司马冏并接掌朝政”的执行中的空间与人员处置”。条目把背景概括为“河间王司马颙与成都王司马颖起兵，长沙王司马乂在洛阳杀齐王司马冏并接掌朝政”发生的政治与区域条件共同作用，并指出该节点改变后续资源、联盟或制度处置的条件。这个节点进入区域社会时，要经由文书、道路、军队、市场、寺院和家庭等不同中介；因此，政治行动的宣布、地方接收和日常生活的改变不能被视作同一时刻完成。

本条依据《晋书·惠帝纪》；《晋书·齐王冏传》；《晋书·长沙厉王乂传》；《晋书·成都王颖传》，证据提示为“核心事项已有既有账本来源；本分解节点的参与者、执行范围和时间先后仍须逐案核验”。这些材料可支持年序、官方决定、战事或特定地点的线索，却难以直接统计所有人的财富、迁徙、信仰和动机。更稳妥的读法是区分诏令与实施、军报与人口、行政称谓与自我认同，并让地方材料的缺环限制解释的范围。

由此可见，该事件不是孤立的年表点，而是资源、信息与权力关系重新配置的一环。不同地区的官员、社群和家庭可能以不同方式承担征发、保护、迁移或宗教活动；现存来源只能照亮其中一部分，不能替沉默者制造统一的经验。

\eventanchor{JH-LEDGER-010-5}{太安元年（302年）·“河间王司马颙与成都王司马颖起兵，长沙王司马乂在洛阳杀齐王司马冏并接掌朝政”的后续制度或军政调整}账本记录西晋在太安元年（302年）发生““河间王司马颙与成都王司马颖起兵，长沙王司马乂在洛阳杀齐王司马冏并接掌朝政”的后续制度或军政调整”。条目把背景概括为核心节点发生后仍须处理其遗留问题，并指出后续影响须在不同地区和材料类型中分别核验。这个节点进入区域社会时，要经由文书、道路、军队、市场、寺院和家庭等不同中介；因此，政治行动的宣布、地方接收和日常生活的改变不能被视作同一时刻完成。

本条依据《晋书·惠帝纪》；《晋书·齐王冏传》；《晋书·长沙厉王乂传》；《晋书·成都王颖传》，证据提示为“核心事项已有既有账本来源；本分解节点的参与者、执行范围和时间先后仍须逐案核验”。这些材料可支持年序、官方决定、战事或特定地点的线索，却难以直接统计所有人的财富、迁徙、信仰和动机。更稳妥的读法是区分诏令与实施、军报与人口、行政称谓与自我认同，并让地方材料的缺环限制解释的范围。

由此可见，该事件不是孤立的年表点，而是资源、信息与权力关系重新配置的一环。不同地区的官员、社群和家庭可能以不同方式承担征发、保护、迁移或宗教活动；现存来源只能照亮其中一部分，不能替沉默者制造统一的经验。

\eventanchor{JH-LEDGER-010-6}{太安元年（302年）·“河间王司马颙与成都王司马颖起兵，长沙王司马乂在洛阳杀齐王司马冏并接掌朝政”的异源材料与影响范围的复核}账本记录西晋在太安元年（302年）发生““河间王司马颙与成都王司马颖起兵，长沙王司马乂在洛阳杀齐王司马冏并接掌朝政”的异源材料与影响范围的复核”。条目把背景概括为核心节点发生后仍须处理其遗留问题，并指出后续影响须在不同地区和材料类型中分别核验。这个节点进入区域社会时，要经由文书、道路、军队、市场、寺院和家庭等不同中介；因此，政治行动的宣布、地方接收和日常生活的改变不能被视作同一时刻完成。

本条依据《晋书·惠帝纪》；《晋书·齐王冏传》；《晋书·长沙厉王乂传》；《晋书·成都王颖传》，证据提示为“核心事项已有既有账本来源；本分解节点的参与者、执行范围和时间先后仍须逐案核验”。这些材料可支持年序、官方决定、战事或特定地点的线索，却难以直接统计所有人的财富、迁徙、信仰和动机。更稳妥的读法是区分诏令与实施、军报与人口、行政称谓与自我认同，并让地方材料的缺环限制解释的范围。

由此可见，该事件不是孤立的年表点，而是资源、信息与权力关系重新配置的一环。不同地区的官员、社群和家庭可能以不同方式承担征发、保护、迁移或宗教活动；现存来源只能照亮其中一部分，不能替沉默者制造统一的经验。

\eventanchor{JH-LEDGER-011-1}{太安二年（303年）·“长沙王司马乂据洛阳抵抗河间王司马颙和成都王司马颖联军，京师陷入长期战斗”的前期动员与资源准备}账本记录西晋在太安二年（303年）发生““长沙王司马乂据洛阳抵抗河间王司马颙和成都王司马颖联军，京师陷入长期战斗”的前期动员与资源准备”。条目把背景概括为当时的权力、资源与信息约束推动相关过程展开，并指出这一过程为“长沙王司马乂据洛阳抵抗河间王司马颙和成都王司马颖联军，京师陷入长期战斗”提供可观察的条件。这个节点进入区域社会时，要经由文书、道路、军队、市场、寺院和家庭等不同中介；因此，政治行动的宣布、地方接收和日常生活的改变不能被视作同一时刻完成。

本条依据《晋书·惠帝纪》；《晋书·长沙厉王乂传》；《晋书·河间平王颙传》；《资治通鉴·晋纪》，证据提示为“核心事项已有既有账本来源；本分解节点的参与者、执行范围和时间先后仍须逐案核验”。这些材料可支持年序、官方决定、战事或特定地点的线索，却难以直接统计所有人的财富、迁徙、信仰和动机。更稳妥的读法是区分诏令与实施、军报与人口、行政称谓与自我认同，并让地方材料的缺环限制解释的范围。

由此可见，该事件不是孤立的年表点，而是资源、信息与权力关系重新配置的一环。不同地区的官员、社群和家庭可能以不同方式承担征发、保护、迁移或宗教活动；现存来源只能照亮其中一部分，不能替沉默者制造统一的经验。

\eventanchor{JH-LEDGER-011-2}{太安二年（303年）·“长沙王司马乂据洛阳抵抗河间王司马颙和成都王司马颖联军，京师陷入长期战斗”的命令、联盟或地方响应}账本记录西晋在太安二年（303年）发生““长沙王司马乂据洛阳抵抗河间王司马颙和成都王司马颖联军，京师陷入长期战斗”的命令、联盟或地方响应”。条目把背景概括为当时的权力、资源与信息约束推动相关过程展开，并指出这一过程为“长沙王司马乂据洛阳抵抗河间王司马颙和成都王司马颖联军，京师陷入长期战斗”提供可观察的条件。这个节点进入区域社会时，要经由文书、道路、军队、市场、寺院和家庭等不同中介；因此，政治行动的宣布、地方接收和日常生活的改变不能被视作同一时刻完成。

本条依据《晋书·惠帝纪》；《晋书·长沙厉王乂传》；《晋书·河间平王颙传》；《资治通鉴·晋纪》，证据提示为“核心事项已有既有账本来源；本分解节点的参与者、执行范围和时间先后仍须逐案核验”。这些材料可支持年序、官方决定、战事或特定地点的线索，却难以直接统计所有人的财富、迁徙、信仰和动机。更稳妥的读法是区分诏令与实施、军报与人口、行政称谓与自我认同，并让地方材料的缺环限制解释的范围。

由此可见，该事件不是孤立的年表点，而是资源、信息与权力关系重新配置的一环。不同地区的官员、社群和家庭可能以不同方式承担征发、保护、迁移或宗教活动；现存来源只能照亮其中一部分，不能替沉默者制造统一的经验。

\eventanchor{JH-LEDGER-011-3}{太安二年（303年）·长沙王司马乂据洛阳抵抗河间王司马颙和成都王司马颖联军，京师陷入长期战斗}账本记录西晋在太安二年（303年）发生“长沙王司马乂据洛阳抵抗河间王司马颙和成都王司马颖联军，京师陷入长期战斗”。条目把背景概括为“长沙王司马乂据洛阳抵抗河间王司马颙和成都王司马颖联军，京师陷入长期战斗”发生的政治与区域条件共同作用，并指出该节点改变后续资源、联盟或制度处置的条件。这个节点进入区域社会时，要经由文书、道路、军队、市场、寺院和家庭等不同中介；因此，政治行动的宣布、地方接收和日常生活的改变不能被视作同一时刻完成。

本条依据《晋书·惠帝纪》；《晋书·长沙厉王乂传》；《晋书·河间平王颙传》；《资治通鉴·晋纪》，证据提示为“核心事项已有既有账本来源；本分解节点的参与者、执行范围和时间先后仍须逐案核验”。这些材料可支持年序、官方决定、战事或特定地点的线索，却难以直接统计所有人的财富、迁徙、信仰和动机。更稳妥的读法是区分诏令与实施、军报与人口、行政称谓与自我认同，并让地方材料的缺环限制解释的范围。

由此可见，该事件不是孤立的年表点，而是资源、信息与权力关系重新配置的一环。不同地区的官员、社群和家庭可能以不同方式承担征发、保护、迁移或宗教活动；现存来源只能照亮其中一部分，不能替沉默者制造统一的经验。

\eventanchor{JH-LEDGER-011-4}{太安二年（303年）·“长沙王司马乂据洛阳抵抗河间王司马颙和成都王司马颖联军，京师陷入长期战斗”的执行中的空间与人员处置}账本记录西晋在太安二年（303年）发生““长沙王司马乂据洛阳抵抗河间王司马颙和成都王司马颖联军，京师陷入长期战斗”的执行中的空间与人员处置”。条目把背景概括为“长沙王司马乂据洛阳抵抗河间王司马颙和成都王司马颖联军，京师陷入长期战斗”发生的政治与区域条件共同作用，并指出该节点改变后续资源、联盟或制度处置的条件。这个节点进入区域社会时，要经由文书、道路、军队、市场、寺院和家庭等不同中介；因此，政治行动的宣布、地方接收和日常生活的改变不能被视作同一时刻完成。

本条依据《晋书·惠帝纪》；《晋书·长沙厉王乂传》；《晋书·河间平王颙传》；《资治通鉴·晋纪》，证据提示为“核心事项已有既有账本来源；本分解节点的参与者、执行范围和时间先后仍须逐案核验”。这些材料可支持年序、官方决定、战事或特定地点的线索，却难以直接统计所有人的财富、迁徙、信仰和动机。更稳妥的读法是区分诏令与实施、军报与人口、行政称谓与自我认同，并让地方材料的缺环限制解释的范围。

由此可见，该事件不是孤立的年表点，而是资源、信息与权力关系重新配置的一环。不同地区的官员、社群和家庭可能以不同方式承担征发、保护、迁移或宗教活动；现存来源只能照亮其中一部分，不能替沉默者制造统一的经验。

\eventanchor{JH-LEDGER-011-5}{太安二年（303年）·“长沙王司马乂据洛阳抵抗河间王司马颙和成都王司马颖联军，京师陷入长期战斗”的后续制度或军政调整}账本记录西晋在太安二年（303年）发生““长沙王司马乂据洛阳抵抗河间王司马颙和成都王司马颖联军，京师陷入长期战斗”的后续制度或军政调整”。条目把背景概括为核心节点发生后仍须处理其遗留问题，并指出后续影响须在不同地区和材料类型中分别核验。这个节点进入区域社会时，要经由文书、道路、军队、市场、寺院和家庭等不同中介；因此，政治行动的宣布、地方接收和日常生活的改变不能被视作同一时刻完成。

本条依据《晋书·惠帝纪》；《晋书·长沙厉王乂传》；《晋书·河间平王颙传》；《资治通鉴·晋纪》，证据提示为“核心事项已有既有账本来源；本分解节点的参与者、执行范围和时间先后仍须逐案核验”。这些材料可支持年序、官方决定、战事或特定地点的线索，却难以直接统计所有人的财富、迁徙、信仰和动机。更稳妥的读法是区分诏令与实施、军报与人口、行政称谓与自我认同，并让地方材料的缺环限制解释的范围。

由此可见，该事件不是孤立的年表点，而是资源、信息与权力关系重新配置的一环。不同地区的官员、社群和家庭可能以不同方式承担征发、保护、迁移或宗教活动；现存来源只能照亮其中一部分，不能替沉默者制造统一的经验。

\eventanchor{JH-LEDGER-011-6}{太安二年（303年）·“长沙王司马乂据洛阳抵抗河间王司马颙和成都王司马颖联军，京师陷入长期战斗”的异源材料与影响范围的复核}账本记录西晋在太安二年（303年）发生““长沙王司马乂据洛阳抵抗河间王司马颙和成都王司马颖联军，京师陷入长期战斗”的异源材料与影响范围的复核”。条目把背景概括为核心节点发生后仍须处理其遗留问题，并指出后续影响须在不同地区和材料类型中分别核验。这个节点进入区域社会时，要经由文书、道路、军队、市场、寺院和家庭等不同中介；因此，政治行动的宣布、地方接收和日常生活的改变不能被视作同一时刻完成。

本条依据《晋书·惠帝纪》；《晋书·长沙厉王乂传》；《晋书·河间平王颙传》；《资治通鉴·晋纪》，证据提示为“核心事项已有既有账本来源；本分解节点的参与者、执行范围和时间先后仍须逐案核验”。这些材料可支持年序、官方决定、战事或特定地点的线索，却难以直接统计所有人的财富、迁徙、信仰和动机。更稳妥的读法是区分诏令与实施、军报与人口、行政称谓与自我认同，并让地方材料的缺环限制解释的范围。

由此可见，该事件不是孤立的年表点，而是资源、信息与权力关系重新配置的一环。不同地区的官员、社群和家庭可能以不同方式承担征发、保护、迁移或宗教活动；现存来源只能照亮其中一部分，不能替沉默者制造统一的经验。

\eventanchor{JH-LEDGER-012-1}{太安二年（303年）·“巴氐首领李雄攻取成都，益州的西晋统治被成汉政权取代”的前期动员与资源准备}账本记录成汉与益州西晋守军在太安二年（303年）发生““巴氐首领李雄攻取成都，益州的西晋统治被成汉政权取代”的前期动员与资源准备”。条目把背景概括为当时的权力、资源与信息约束推动相关过程展开，并指出这一过程为“巴氐首领李雄攻取成都，益州的西晋统治被成汉政权取代”提供可观察的条件。这个节点进入区域社会时，要经由文书、道路、军队、市场、寺院和家庭等不同中介；因此，政治行动的宣布、地方接收和日常生活的改变不能被视作同一时刻完成。

本条依据《晋书·李雄载记》；《华阳国志·李特雄期寿势志》；《资治通鉴·晋纪》，证据提示为“核心事项已有既有账本来源；本分解节点的参与者、执行范围和时间先后仍须逐案核验”。这些材料可支持年序、官方决定、战事或特定地点的线索，却难以直接统计所有人的财富、迁徙、信仰和动机。更稳妥的读法是区分诏令与实施、军报与人口、行政称谓与自我认同，并让地方材料的缺环限制解释的范围。

由此可见，该事件不是孤立的年表点，而是资源、信息与权力关系重新配置的一环。不同地区的官员、社群和家庭可能以不同方式承担征发、保护、迁移或宗教活动；现存来源只能照亮其中一部分，不能替沉默者制造统一的经验。

\eventanchor{JH-LEDGER-012-2}{太安二年（303年）·“巴氐首领李雄攻取成都，益州的西晋统治被成汉政权取代”的命令、联盟或地方响应}账本记录成汉与益州西晋守军在太安二年（303年）发生““巴氐首领李雄攻取成都，益州的西晋统治被成汉政权取代”的命令、联盟或地方响应”。条目把背景概括为当时的权力、资源与信息约束推动相关过程展开，并指出这一过程为“巴氐首领李雄攻取成都，益州的西晋统治被成汉政权取代”提供可观察的条件。这个节点进入区域社会时，要经由文书、道路、军队、市场、寺院和家庭等不同中介；因此，政治行动的宣布、地方接收和日常生活的改变不能被视作同一时刻完成。

本条依据《晋书·李雄载记》；《华阳国志·李特雄期寿势志》；《资治通鉴·晋纪》，证据提示为“核心事项已有既有账本来源；本分解节点的参与者、执行范围和时间先后仍须逐案核验”。这些材料可支持年序、官方决定、战事或特定地点的线索，却难以直接统计所有人的财富、迁徙、信仰和动机。更稳妥的读法是区分诏令与实施、军报与人口、行政称谓与自我认同，并让地方材料的缺环限制解释的范围。

由此可见，该事件不是孤立的年表点，而是资源、信息与权力关系重新配置的一环。不同地区的官员、社群和家庭可能以不同方式承担征发、保护、迁移或宗教活动；现存来源只能照亮其中一部分，不能替沉默者制造统一的经验。

\eventanchor{JH-LEDGER-012-3}{太安二年（303年）·巴氐首领李雄攻取成都，益州的西晋统治被成汉政权取代}账本记录成汉与益州西晋守军在太安二年（303年）发生“巴氐首领李雄攻取成都，益州的西晋统治被成汉政权取代”。条目把背景概括为“巴氐首领李雄攻取成都，益州的西晋统治被成汉政权取代”发生的政治与区域条件共同作用，并指出该节点改变后续资源、联盟或制度处置的条件。这个节点进入区域社会时，要经由文书、道路、军队、市场、寺院和家庭等不同中介；因此，政治行动的宣布、地方接收和日常生活的改变不能被视作同一时刻完成。

本条依据《晋书·李雄载记》；《华阳国志·李特雄期寿势志》；《资治通鉴·晋纪》，证据提示为“核心事项已有既有账本来源；本分解节点的参与者、执行范围和时间先后仍须逐案核验”。这些材料可支持年序、官方决定、战事或特定地点的线索，却难以直接统计所有人的财富、迁徙、信仰和动机。更稳妥的读法是区分诏令与实施、军报与人口、行政称谓与自我认同，并让地方材料的缺环限制解释的范围。

由此可见，该事件不是孤立的年表点，而是资源、信息与权力关系重新配置的一环。不同地区的官员、社群和家庭可能以不同方式承担征发、保护、迁移或宗教活动；现存来源只能照亮其中一部分，不能替沉默者制造统一的经验。

\eventanchor{JH-LEDGER-012-4}{太安二年（303年）·“巴氐首领李雄攻取成都，益州的西晋统治被成汉政权取代”的执行中的空间与人员处置}账本记录成汉与益州西晋守军在太安二年（303年）发生““巴氐首领李雄攻取成都，益州的西晋统治被成汉政权取代”的执行中的空间与人员处置”。条目把背景概括为“巴氐首领李雄攻取成都，益州的西晋统治被成汉政权取代”发生的政治与区域条件共同作用，并指出该节点改变后续资源、联盟或制度处置的条件。这个节点进入区域社会时，要经由文书、道路、军队、市场、寺院和家庭等不同中介；因此，政治行动的宣布、地方接收和日常生活的改变不能被视作同一时刻完成。

本条依据《晋书·李雄载记》；《华阳国志·李特雄期寿势志》；《资治通鉴·晋纪》，证据提示为“核心事项已有既有账本来源；本分解节点的参与者、执行范围和时间先后仍须逐案核验”。这些材料可支持年序、官方决定、战事或特定地点的线索，却难以直接统计所有人的财富、迁徙、信仰和动机。更稳妥的读法是区分诏令与实施、军报与人口、行政称谓与自我认同，并让地方材料的缺环限制解释的范围。

由此可见，该事件不是孤立的年表点，而是资源、信息与权力关系重新配置的一环。不同地区的官员、社群和家庭可能以不同方式承担征发、保护、迁移或宗教活动；现存来源只能照亮其中一部分，不能替沉默者制造统一的经验。

\eventanchor{JH-LEDGER-012-5}{太安二年（303年）·“巴氐首领李雄攻取成都，益州的西晋统治被成汉政权取代”的后续制度或军政调整}账本记录成汉与益州西晋守军在太安二年（303年）发生““巴氐首领李雄攻取成都，益州的西晋统治被成汉政权取代”的后续制度或军政调整”。条目把背景概括为核心节点发生后仍须处理其遗留问题，并指出后续影响须在不同地区和材料类型中分别核验。这个节点进入区域社会时，要经由文书、道路、军队、市场、寺院和家庭等不同中介；因此，政治行动的宣布、地方接收和日常生活的改变不能被视作同一时刻完成。

本条依据《晋书·李雄载记》；《华阳国志·李特雄期寿势志》；《资治通鉴·晋纪》，证据提示为“核心事项已有既有账本来源；本分解节点的参与者、执行范围和时间先后仍须逐案核验”。这些材料可支持年序、官方决定、战事或特定地点的线索，却难以直接统计所有人的财富、迁徙、信仰和动机。更稳妥的读法是区分诏令与实施、军报与人口、行政称谓与自我认同，并让地方材料的缺环限制解释的范围。

由此可见，该事件不是孤立的年表点，而是资源、信息与权力关系重新配置的一环。不同地区的官员、社群和家庭可能以不同方式承担征发、保护、迁移或宗教活动；现存来源只能照亮其中一部分，不能替沉默者制造统一的经验。

\eventanchor{JH-LEDGER-012-6}{太安二年（303年）·“巴氐首领李雄攻取成都，益州的西晋统治被成汉政权取代”的异源材料与影响范围的复核}账本记录成汉与益州西晋守军在太安二年（303年）发生““巴氐首领李雄攻取成都，益州的西晋统治被成汉政权取代”的异源材料与影响范围的复核”。条目把背景概括为核心节点发生后仍须处理其遗留问题，并指出后续影响须在不同地区和材料类型中分别核验。这个节点进入区域社会时，要经由文书、道路、军队、市场、寺院和家庭等不同中介；因此，政治行动的宣布、地方接收和日常生活的改变不能被视作同一时刻完成。

本条依据《晋书·李雄载记》；《华阳国志·李特雄期寿势志》；《资治通鉴·晋纪》，证据提示为“核心事项已有既有账本来源；本分解节点的参与者、执行范围和时间先后仍须逐案核验”。这些材料可支持年序、官方决定、战事或特定地点的线索，却难以直接统计所有人的财富、迁徙、信仰和动机。更稳妥的读法是区分诏令与实施、军报与人口、行政称谓与自我认同，并让地方材料的缺环限制解释的范围。

由此可见，该事件不是孤立的年表点，而是资源、信息与权力关系重新配置的一环。不同地区的官员、社群和家庭可能以不同方式承担征发、保护、迁移或宗教活动；现存来源只能照亮其中一部分，不能替沉默者制造统一的经验。

\eventanchor{JH-LEDGER-013-1}{永安元年（304年）·“司马乂被东海王司马越等交给河间王司马颙军并遭杀害，洛阳政权落入外军支配”的前期动员与资源准备}账本记录西晋在永安元年（304年）发生““司马乂被东海王司马越等交给河间王司马颙军并遭杀害，洛阳政权落入外军支配”的前期动员与资源准备”。条目把背景概括为当时的权力、资源与信息约束推动相关过程展开，并指出这一过程为“司马乂被东海王司马越等交给河间王司马颙军并遭杀害，洛阳政权落入外军支配”提供可观察的条件。这个节点进入区域社会时，要经由文书、道路、军队、市场、寺院和家庭等不同中介；因此，政治行动的宣布、地方接收和日常生活的改变不能被视作同一时刻完成。

本条依据《晋书·惠帝纪》；《晋书·长沙厉王乂传》；《晋书·东海孝献王越传》，证据提示为“核心事项已有既有账本来源；本分解节点的参与者、执行范围和时间先后仍须逐案核验”。这些材料可支持年序、官方决定、战事或特定地点的线索，却难以直接统计所有人的财富、迁徙、信仰和动机。更稳妥的读法是区分诏令与实施、军报与人口、行政称谓与自我认同，并让地方材料的缺环限制解释的范围。

由此可见，该事件不是孤立的年表点，而是资源、信息与权力关系重新配置的一环。不同地区的官员、社群和家庭可能以不同方式承担征发、保护、迁移或宗教活动；现存来源只能照亮其中一部分，不能替沉默者制造统一的经验。

\eventanchor{JH-LEDGER-013-2}{永安元年（304年）·“司马乂被东海王司马越等交给河间王司马颙军并遭杀害，洛阳政权落入外军支配”的命令、联盟或地方响应}账本记录西晋在永安元年（304年）发生““司马乂被东海王司马越等交给河间王司马颙军并遭杀害，洛阳政权落入外军支配”的命令、联盟或地方响应”。条目把背景概括为当时的权力、资源与信息约束推动相关过程展开，并指出这一过程为“司马乂被东海王司马越等交给河间王司马颙军并遭杀害，洛阳政权落入外军支配”提供可观察的条件。这个节点进入区域社会时，要经由文书、道路、军队、市场、寺院和家庭等不同中介；因此，政治行动的宣布、地方接收和日常生活的改变不能被视作同一时刻完成。

本条依据《晋书·惠帝纪》；《晋书·长沙厉王乂传》；《晋书·东海孝献王越传》，证据提示为“核心事项已有既有账本来源；本分解节点的参与者、执行范围和时间先后仍须逐案核验”。这些材料可支持年序、官方决定、战事或特定地点的线索，却难以直接统计所有人的财富、迁徙、信仰和动机。更稳妥的读法是区分诏令与实施、军报与人口、行政称谓与自我认同，并让地方材料的缺环限制解释的范围。

由此可见，该事件不是孤立的年表点，而是资源、信息与权力关系重新配置的一环。不同地区的官员、社群和家庭可能以不同方式承担征发、保护、迁移或宗教活动；现存来源只能照亮其中一部分，不能替沉默者制造统一的经验。

\subsection{社会关系与宗教空间}

\eventanchor{JH-LEDGER-013-3}{永安元年（304年）·司马乂被东海王司马越等交给河间王司马颙军并遭杀害，洛阳政权落入外军支配}账本记录西晋在永安元年（304年）发生“司马乂被东海王司马越等交给河间王司马颙军并遭杀害，洛阳政权落入外军支配”。条目把背景概括为“司马乂被东海王司马越等交给河间王司马颙军并遭杀害，洛阳政权落入外军支配”发生的政治与区域条件共同作用，并指出该节点改变后续资源、联盟或制度处置的条件。这个节点进入区域社会时，要经由文书、道路、军队、市场、寺院和家庭等不同中介；因此，政治行动的宣布、地方接收和日常生活的改变不能被视作同一时刻完成。

本条依据《晋书·惠帝纪》；《晋书·长沙厉王乂传》；《晋书·东海孝献王越传》，证据提示为“核心事项已有既有账本来源；本分解节点的参与者、执行范围和时间先后仍须逐案核验”。这些材料可支持年序、官方决定、战事或特定地点的线索，却难以直接统计所有人的财富、迁徙、信仰和动机。更稳妥的读法是区分诏令与实施、军报与人口、行政称谓与自我认同，并让地方材料的缺环限制解释的范围。

由此可见，该事件不是孤立的年表点，而是资源、信息与权力关系重新配置的一环。不同地区的官员、社群和家庭可能以不同方式承担征发、保护、迁移或宗教活动；现存来源只能照亮其中一部分，不能替沉默者制造统一的经验。

\eventanchor{JH-LEDGER-013-4}{永安元年（304年）·“司马乂被东海王司马越等交给河间王司马颙军并遭杀害，洛阳政权落入外军支配”的执行中的空间与人员处置}账本记录西晋在永安元年（304年）发生““司马乂被东海王司马越等交给河间王司马颙军并遭杀害，洛阳政权落入外军支配”的执行中的空间与人员处置”。条目把背景概括为“司马乂被东海王司马越等交给河间王司马颙军并遭杀害，洛阳政权落入外军支配”发生的政治与区域条件共同作用，并指出该节点改变后续资源、联盟或制度处置的条件。这个节点进入区域社会时，要经由文书、道路、军队、市场、寺院和家庭等不同中介；因此，政治行动的宣布、地方接收和日常生活的改变不能被视作同一时刻完成。

本条依据《晋书·惠帝纪》；《晋书·长沙厉王乂传》；《晋书·东海孝献王越传》，证据提示为“核心事项已有既有账本来源；本分解节点的参与者、执行范围和时间先后仍须逐案核验”。这些材料可支持年序、官方决定、战事或特定地点的线索，却难以直接统计所有人的财富、迁徙、信仰和动机。更稳妥的读法是区分诏令与实施、军报与人口、行政称谓与自我认同，并让地方材料的缺环限制解释的范围。

由此可见，该事件不是孤立的年表点，而是资源、信息与权力关系重新配置的一环。不同地区的官员、社群和家庭可能以不同方式承担征发、保护、迁移或宗教活动；现存来源只能照亮其中一部分，不能替沉默者制造统一的经验。

\eventanchor{JH-LEDGER-013-5}{永安元年（304年）·“司马乂被东海王司马越等交给河间王司马颙军并遭杀害，洛阳政权落入外军支配”的后续制度或军政调整}账本记录西晋在永安元年（304年）发生““司马乂被东海王司马越等交给河间王司马颙军并遭杀害，洛阳政权落入外军支配”的后续制度或军政调整”。条目把背景概括为核心节点发生后仍须处理其遗留问题，并指出后续影响须在不同地区和材料类型中分别核验。这个节点进入区域社会时，要经由文书、道路、军队、市场、寺院和家庭等不同中介；因此，政治行动的宣布、地方接收和日常生活的改变不能被视作同一时刻完成。

本条依据《晋书·惠帝纪》；《晋书·长沙厉王乂传》；《晋书·东海孝献王越传》，证据提示为“核心事项已有既有账本来源；本分解节点的参与者、执行范围和时间先后仍须逐案核验”。这些材料可支持年序、官方决定、战事或特定地点的线索，却难以直接统计所有人的财富、迁徙、信仰和动机。更稳妥的读法是区分诏令与实施、军报与人口、行政称谓与自我认同，并让地方材料的缺环限制解释的范围。

由此可见，该事件不是孤立的年表点，而是资源、信息与权力关系重新配置的一环。不同地区的官员、社群和家庭可能以不同方式承担征发、保护、迁移或宗教活动；现存来源只能照亮其中一部分，不能替沉默者制造统一的经验。

\eventanchor{JH-LEDGER-013-6}{永安元年（304年）·“司马乂被东海王司马越等交给河间王司马颙军并遭杀害，洛阳政权落入外军支配”的异源材料与影响范围的复核}账本记录西晋在永安元年（304年）发生““司马乂被东海王司马越等交给河间王司马颙军并遭杀害，洛阳政权落入外军支配”的异源材料与影响范围的复核”。条目把背景概括为核心节点发生后仍须处理其遗留问题，并指出后续影响须在不同地区和材料类型中分别核验。这个节点进入区域社会时，要经由文书、道路、军队、市场、寺院和家庭等不同中介；因此，政治行动的宣布、地方接收和日常生活的改变不能被视作同一时刻完成。

本条依据《晋书·惠帝纪》；《晋书·长沙厉王乂传》；《晋书·东海孝献王越传》，证据提示为“核心事项已有既有账本来源；本分解节点的参与者、执行范围和时间先后仍须逐案核验”。这些材料可支持年序、官方决定、战事或特定地点的线索，却难以直接统计所有人的财富、迁徙、信仰和动机。更稳妥的读法是区分诏令与实施、军报与人口、行政称谓与自我认同，并让地方材料的缺环限制解释的范围。

由此可见，该事件不是孤立的年表点，而是资源、信息与权力关系重新配置的一环。不同地区的官员、社群和家庭可能以不同方式承担征发、保护、迁移或宗教活动；现存来源只能照亮其中一部分，不能替沉默者制造统一的经验。

\eventanchor{JH-LEDGER-014-1}{永兴元年（304年）·“刘渊在左国城起兵，自称汉王，汇聚匈奴五部及并州反晋力量”的前期动员与资源准备}账本记录汉赵与并州诸部在永兴元年（304年）发生““刘渊在左国城起兵，自称汉王，汇聚匈奴五部及并州反晋力量”的前期动员与资源准备”。条目把背景概括为当时的权力、资源与信息约束推动相关过程展开，并指出这一过程为“刘渊在左国城起兵，自称汉王，汇聚匈奴五部及并州反晋力量”提供可观察的条件。这个节点进入区域社会时，要经由文书、道路、军队、市场、寺院和家庭等不同中介；因此，政治行动的宣布、地方接收和日常生活的改变不能被视作同一时刻完成。

本条依据《晋书·刘元海载记》；《资治通鉴·晋纪》；《十六国春秋》辑佚，证据提示为“核心事项已有既有账本来源；本分解节点的参与者、执行范围和时间先后仍须逐案核验”。这些材料可支持年序、官方决定、战事或特定地点的线索，却难以直接统计所有人的财富、迁徙、信仰和动机。更稳妥的读法是区分诏令与实施、军报与人口、行政称谓与自我认同，并让地方材料的缺环限制解释的范围。

由此可见，该事件不是孤立的年表点，而是资源、信息与权力关系重新配置的一环。不同地区的官员、社群和家庭可能以不同方式承担征发、保护、迁移或宗教活动；现存来源只能照亮其中一部分，不能替沉默者制造统一的经验。

\eventanchor{JH-LEDGER-014-2}{永兴元年（304年）·“刘渊在左国城起兵，自称汉王，汇聚匈奴五部及并州反晋力量”的命令、联盟或地方响应}账本记录汉赵与并州诸部在永兴元年（304年）发生““刘渊在左国城起兵，自称汉王，汇聚匈奴五部及并州反晋力量”的命令、联盟或地方响应”。条目把背景概括为当时的权力、资源与信息约束推动相关过程展开，并指出这一过程为“刘渊在左国城起兵，自称汉王，汇聚匈奴五部及并州反晋力量”提供可观察的条件。这个节点进入区域社会时，要经由文书、道路、军队、市场、寺院和家庭等不同中介；因此，政治行动的宣布、地方接收和日常生活的改变不能被视作同一时刻完成。

本条依据《晋书·刘元海载记》；《资治通鉴·晋纪》；《十六国春秋》辑佚，证据提示为“核心事项已有既有账本来源；本分解节点的参与者、执行范围和时间先后仍须逐案核验”。这些材料可支持年序、官方决定、战事或特定地点的线索，却难以直接统计所有人的财富、迁徙、信仰和动机。更稳妥的读法是区分诏令与实施、军报与人口、行政称谓与自我认同，并让地方材料的缺环限制解释的范围。

由此可见，该事件不是孤立的年表点，而是资源、信息与权力关系重新配置的一环。不同地区的官员、社群和家庭可能以不同方式承担征发、保护、迁移或宗教活动；现存来源只能照亮其中一部分，不能替沉默者制造统一的经验。

\eventanchor{JH-LEDGER-014-3}{永兴元年（304年）·刘渊在左国城起兵，自称汉王，汇聚匈奴五部及并州反晋力量}账本记录汉赵与并州诸部在永兴元年（304年）发生“刘渊在左国城起兵，自称汉王，汇聚匈奴五部及并州反晋力量”。条目把背景概括为“刘渊在左国城起兵，自称汉王，汇聚匈奴五部及并州反晋力量”发生的政治与区域条件共同作用，并指出该节点改变后续资源、联盟或制度处置的条件。这个节点进入区域社会时，要经由文书、道路、军队、市场、寺院和家庭等不同中介；因此，政治行动的宣布、地方接收和日常生活的改变不能被视作同一时刻完成。

本条依据《晋书·刘元海载记》；《资治通鉴·晋纪》；《十六国春秋》辑佚，证据提示为“核心事项已有既有账本来源；本分解节点的参与者、执行范围和时间先后仍须逐案核验”。这些材料可支持年序、官方决定、战事或特定地点的线索，却难以直接统计所有人的财富、迁徙、信仰和动机。更稳妥的读法是区分诏令与实施、军报与人口、行政称谓与自我认同，并让地方材料的缺环限制解释的范围。

由此可见，该事件不是孤立的年表点，而是资源、信息与权力关系重新配置的一环。不同地区的官员、社群和家庭可能以不同方式承担征发、保护、迁移或宗教活动；现存来源只能照亮其中一部分，不能替沉默者制造统一的经验。

\eventanchor{JH-LEDGER-014-4}{永兴元年（304年）·“刘渊在左国城起兵，自称汉王，汇聚匈奴五部及并州反晋力量”的执行中的空间与人员处置}账本记录汉赵与并州诸部在永兴元年（304年）发生““刘渊在左国城起兵，自称汉王，汇聚匈奴五部及并州反晋力量”的执行中的空间与人员处置”。条目把背景概括为“刘渊在左国城起兵，自称汉王，汇聚匈奴五部及并州反晋力量”发生的政治与区域条件共同作用，并指出该节点改变后续资源、联盟或制度处置的条件。这个节点进入区域社会时，要经由文书、道路、军队、市场、寺院和家庭等不同中介；因此，政治行动的宣布、地方接收和日常生活的改变不能被视作同一时刻完成。

本条依据《晋书·刘元海载记》；《资治通鉴·晋纪》；《十六国春秋》辑佚，证据提示为“核心事项已有既有账本来源；本分解节点的参与者、执行范围和时间先后仍须逐案核验”。这些材料可支持年序、官方决定、战事或特定地点的线索，却难以直接统计所有人的财富、迁徙、信仰和动机。更稳妥的读法是区分诏令与实施、军报与人口、行政称谓与自我认同，并让地方材料的缺环限制解释的范围。

由此可见，该事件不是孤立的年表点，而是资源、信息与权力关系重新配置的一环。不同地区的官员、社群和家庭可能以不同方式承担征发、保护、迁移或宗教活动；现存来源只能照亮其中一部分，不能替沉默者制造统一的经验。

\eventanchor{JH-LEDGER-014-5}{永兴元年（304年）·“刘渊在左国城起兵，自称汉王，汇聚匈奴五部及并州反晋力量”的后续制度或军政调整}账本记录汉赵与并州诸部在永兴元年（304年）发生““刘渊在左国城起兵，自称汉王，汇聚匈奴五部及并州反晋力量”的后续制度或军政调整”。条目把背景概括为核心节点发生后仍须处理其遗留问题，并指出后续影响须在不同地区和材料类型中分别核验。这个节点进入区域社会时，要经由文书、道路、军队、市场、寺院和家庭等不同中介；因此，政治行动的宣布、地方接收和日常生活的改变不能被视作同一时刻完成。

本条依据《晋书·刘元海载记》；《资治通鉴·晋纪》；《十六国春秋》辑佚，证据提示为“核心事项已有既有账本来源；本分解节点的参与者、执行范围和时间先后仍须逐案核验”。这些材料可支持年序、官方决定、战事或特定地点的线索，却难以直接统计所有人的财富、迁徙、信仰和动机。更稳妥的读法是区分诏令与实施、军报与人口、行政称谓与自我认同，并让地方材料的缺环限制解释的范围。

由此可见，该事件不是孤立的年表点，而是资源、信息与权力关系重新配置的一环。不同地区的官员、社群和家庭可能以不同方式承担征发、保护、迁移或宗教活动；现存来源只能照亮其中一部分，不能替沉默者制造统一的经验。

\eventanchor{JH-LEDGER-014-6}{永兴元年（304年）·“刘渊在左国城起兵，自称汉王，汇聚匈奴五部及并州反晋力量”的异源材料与影响范围的复核}账本记录汉赵与并州诸部在永兴元年（304年）发生““刘渊在左国城起兵，自称汉王，汇聚匈奴五部及并州反晋力量”的异源材料与影响范围的复核”。条目把背景概括为核心节点发生后仍须处理其遗留问题，并指出后续影响须在不同地区和材料类型中分别核验。这个节点进入区域社会时，要经由文书、道路、军队、市场、寺院和家庭等不同中介；因此，政治行动的宣布、地方接收和日常生活的改变不能被视作同一时刻完成。

本条依据《晋书·刘元海载记》；《资治通鉴·晋纪》；《十六国春秋》辑佚，证据提示为“核心事项已有既有账本来源；本分解节点的参与者、执行范围和时间先后仍须逐案核验”。这些材料可支持年序、官方决定、战事或特定地点的线索，却难以直接统计所有人的财富、迁徙、信仰和动机。更稳妥的读法是区分诏令与实施、军报与人口、行政称谓与自我认同，并让地方材料的缺环限制解释的范围。

由此可见，该事件不是孤立的年表点，而是资源、信息与权力关系重新配置的一环。不同地区的官员、社群和家庭可能以不同方式承担征发、保护、迁移或宗教活动；现存来源只能照亮其中一部分，不能替沉默者制造统一的经验。

\eventanchor{JH-LEDGER-015-1}{永兴元年（304年）·“陈敏据江东起兵并自置官属，试图利用八王之乱切割东南地区”的前期动员与资源准备}账本记录西晋与江东地方集团在永兴元年（304年）发生““陈敏据江东起兵并自置官属，试图利用八王之乱切割东南地区”的前期动员与资源准备”。条目把背景概括为当时的权力、资源与信息约束推动相关过程展开，并指出这一过程为“陈敏据江东起兵并自置官属，试图利用八王之乱切割东南地区”提供可观察的条件。这个节点进入区域社会时，要经由文书、道路、军队、市场、寺院和家庭等不同中介；因此，政治行动的宣布、地方接收和日常生活的改变不能被视作同一时刻完成。

本条依据《晋书·陈敏传》；《资治通鉴·晋纪》；《建康实录》，证据提示为“核心事项已有既有账本来源；本分解节点的参与者、执行范围和时间先后仍须逐案核验”。这些材料可支持年序、官方决定、战事或特定地点的线索，却难以直接统计所有人的财富、迁徙、信仰和动机。更稳妥的读法是区分诏令与实施、军报与人口、行政称谓与自我认同，并让地方材料的缺环限制解释的范围。

由此可见，该事件不是孤立的年表点，而是资源、信息与权力关系重新配置的一环。不同地区的官员、社群和家庭可能以不同方式承担征发、保护、迁移或宗教活动；现存来源只能照亮其中一部分，不能替沉默者制造统一的经验。

\eventanchor{JH-LEDGER-015-2}{永兴元年（304年）·“陈敏据江东起兵并自置官属，试图利用八王之乱切割东南地区”的命令、联盟或地方响应}账本记录西晋与江东地方集团在永兴元年（304年）发生““陈敏据江东起兵并自置官属，试图利用八王之乱切割东南地区”的命令、联盟或地方响应”。条目把背景概括为当时的权力、资源与信息约束推动相关过程展开，并指出这一过程为“陈敏据江东起兵并自置官属，试图利用八王之乱切割东南地区”提供可观察的条件。这个节点进入区域社会时，要经由文书、道路、军队、市场、寺院和家庭等不同中介；因此，政治行动的宣布、地方接收和日常生活的改变不能被视作同一时刻完成。

本条依据《晋书·陈敏传》；《资治通鉴·晋纪》；《建康实录》，证据提示为“核心事项已有既有账本来源；本分解节点的参与者、执行范围和时间先后仍须逐案核验”。这些材料可支持年序、官方决定、战事或特定地点的线索，却难以直接统计所有人的财富、迁徙、信仰和动机。更稳妥的读法是区分诏令与实施、军报与人口、行政称谓与自我认同，并让地方材料的缺环限制解释的范围。

由此可见，该事件不是孤立的年表点，而是资源、信息与权力关系重新配置的一环。不同地区的官员、社群和家庭可能以不同方式承担征发、保护、迁移或宗教活动；现存来源只能照亮其中一部分，不能替沉默者制造统一的经验。

\eventanchor{JH-LEDGER-015-3}{永兴元年（304年）·陈敏据江东起兵并自置官属，试图利用八王之乱切割东南地区}账本记录西晋与江东地方集团在永兴元年（304年）发生“陈敏据江东起兵并自置官属，试图利用八王之乱切割东南地区”。条目把背景概括为“陈敏据江东起兵并自置官属，试图利用八王之乱切割东南地区”发生的政治与区域条件共同作用，并指出该节点改变后续资源、联盟或制度处置的条件。这个节点进入区域社会时，要经由文书、道路、军队、市场、寺院和家庭等不同中介；因此，政治行动的宣布、地方接收和日常生活的改变不能被视作同一时刻完成。

本条依据《晋书·陈敏传》；《资治通鉴·晋纪》；《建康实录》，证据提示为“核心事项已有既有账本来源；本分解节点的参与者、执行范围和时间先后仍须逐案核验”。这些材料可支持年序、官方决定、战事或特定地点的线索，却难以直接统计所有人的财富、迁徙、信仰和动机。更稳妥的读法是区分诏令与实施、军报与人口、行政称谓与自我认同，并让地方材料的缺环限制解释的范围。

由此可见，该事件不是孤立的年表点，而是资源、信息与权力关系重新配置的一环。不同地区的官员、社群和家庭可能以不同方式承担征发、保护、迁移或宗教活动；现存来源只能照亮其中一部分，不能替沉默者制造统一的经验。

\eventanchor{JH-LEDGER-015-4}{永兴元年（304年）·“陈敏据江东起兵并自置官属，试图利用八王之乱切割东南地区”的执行中的空间与人员处置}账本记录西晋与江东地方集团在永兴元年（304年）发生““陈敏据江东起兵并自置官属，试图利用八王之乱切割东南地区”的执行中的空间与人员处置”。条目把背景概括为“陈敏据江东起兵并自置官属，试图利用八王之乱切割东南地区”发生的政治与区域条件共同作用，并指出该节点改变后续资源、联盟或制度处置的条件。这个节点进入区域社会时，要经由文书、道路、军队、市场、寺院和家庭等不同中介；因此，政治行动的宣布、地方接收和日常生活的改变不能被视作同一时刻完成。

本条依据《晋书·陈敏传》；《资治通鉴·晋纪》；《建康实录》，证据提示为“核心事项已有既有账本来源；本分解节点的参与者、执行范围和时间先后仍须逐案核验”。这些材料可支持年序、官方决定、战事或特定地点的线索，却难以直接统计所有人的财富、迁徙、信仰和动机。更稳妥的读法是区分诏令与实施、军报与人口、行政称谓与自我认同，并让地方材料的缺环限制解释的范围。

由此可见，该事件不是孤立的年表点，而是资源、信息与权力关系重新配置的一环。不同地区的官员、社群和家庭可能以不同方式承担征发、保护、迁移或宗教活动；现存来源只能照亮其中一部分，不能替沉默者制造统一的经验。

\eventanchor{JH-LEDGER-015-5}{永兴元年（304年）·“陈敏据江东起兵并自置官属，试图利用八王之乱切割东南地区”的后续制度或军政调整}账本记录西晋与江东地方集团在永兴元年（304年）发生““陈敏据江东起兵并自置官属，试图利用八王之乱切割东南地区”的后续制度或军政调整”。条目把背景概括为核心节点发生后仍须处理其遗留问题，并指出后续影响须在不同地区和材料类型中分别核验。这个节点进入区域社会时，要经由文书、道路、军队、市场、寺院和家庭等不同中介；因此，政治行动的宣布、地方接收和日常生活的改变不能被视作同一时刻完成。

本条依据《晋书·陈敏传》；《资治通鉴·晋纪》；《建康实录》，证据提示为“核心事项已有既有账本来源；本分解节点的参与者、执行范围和时间先后仍须逐案核验”。这些材料可支持年序、官方决定、战事或特定地点的线索，却难以直接统计所有人的财富、迁徙、信仰和动机。更稳妥的读法是区分诏令与实施、军报与人口、行政称谓与自我认同，并让地方材料的缺环限制解释的范围。

由此可见，该事件不是孤立的年表点，而是资源、信息与权力关系重新配置的一环。不同地区的官员、社群和家庭可能以不同方式承担征发、保护、迁移或宗教活动；现存来源只能照亮其中一部分，不能替沉默者制造统一的经验。

\eventanchor{JH-LEDGER-015-6}{永兴元年（304年）·“陈敏据江东起兵并自置官属，试图利用八王之乱切割东南地区”的异源材料与影响范围的复核}账本记录西晋与江东地方集团在永兴元年（304年）发生““陈敏据江东起兵并自置官属，试图利用八王之乱切割东南地区”的异源材料与影响范围的复核”。条目把背景概括为核心节点发生后仍须处理其遗留问题，并指出后续影响须在不同地区和材料类型中分别核验。这个节点进入区域社会时，要经由文书、道路、军队、市场、寺院和家庭等不同中介；因此，政治行动的宣布、地方接收和日常生活的改变不能被视作同一时刻完成。

本条依据《晋书·陈敏传》；《资治通鉴·晋纪》；《建康实录》，证据提示为“核心事项已有既有账本来源；本分解节点的参与者、执行范围和时间先后仍须逐案核验”。这些材料可支持年序、官方决定、战事或特定地点的线索，却难以直接统计所有人的财富、迁徙、信仰和动机。更稳妥的读法是区分诏令与实施、军报与人口、行政称谓与自我认同，并让地方材料的缺环限制解释的范围。

由此可见，该事件不是孤立的年表点，而是资源、信息与权力关系重新配置的一环。不同地区的官员、社群和家庭可能以不同方式承担征发、保护、迁移或宗教活动；现存来源只能照亮其中一部分，不能替沉默者制造统一的经验。

\eventanchor{JH-LEDGER-016-1}{永兴二年（305年）·“东海王司马越以讨河间王司马颙为名起兵，八王之乱扩大为跨州郡的长期战争”的前期动员与资源准备}账本记录西晋在永兴二年（305年）发生““东海王司马越以讨河间王司马颙为名起兵，八王之乱扩大为跨州郡的长期战争”的前期动员与资源准备”。条目把背景概括为当时的权力、资源与信息约束推动相关过程展开，并指出这一过程为“东海王司马越以讨河间王司马颙为名起兵，八王之乱扩大为跨州郡的长期战争”提供可观察的条件。这个节点进入区域社会时，要经由文书、道路、军队、市场、寺院和家庭等不同中介；因此，政治行动的宣布、地方接收和日常生活的改变不能被视作同一时刻完成。

本条依据《晋书·惠帝纪》；《晋书·东海孝献王越传》；《晋书·河间平王颙传》，证据提示为“核心事项已有既有账本来源；本分解节点的参与者、执行范围和时间先后仍须逐案核验”。这些材料可支持年序、官方决定、战事或特定地点的线索，却难以直接统计所有人的财富、迁徙、信仰和动机。更稳妥的读法是区分诏令与实施、军报与人口、行政称谓与自我认同，并让地方材料的缺环限制解释的范围。

由此可见，该事件不是孤立的年表点，而是资源、信息与权力关系重新配置的一环。不同地区的官员、社群和家庭可能以不同方式承担征发、保护、迁移或宗教活动；现存来源只能照亮其中一部分，不能替沉默者制造统一的经验。

\eventanchor{JH-LEDGER-016-2}{永兴二年（305年）·“东海王司马越以讨河间王司马颙为名起兵，八王之乱扩大为跨州郡的长期战争”的命令、联盟或地方响应}账本记录西晋在永兴二年（305年）发生““东海王司马越以讨河间王司马颙为名起兵，八王之乱扩大为跨州郡的长期战争”的命令、联盟或地方响应”。条目把背景概括为当时的权力、资源与信息约束推动相关过程展开，并指出这一过程为“东海王司马越以讨河间王司马颙为名起兵，八王之乱扩大为跨州郡的长期战争”提供可观察的条件。这个节点进入区域社会时，要经由文书、道路、军队、市场、寺院和家庭等不同中介；因此，政治行动的宣布、地方接收和日常生活的改变不能被视作同一时刻完成。

本条依据《晋书·惠帝纪》；《晋书·东海孝献王越传》；《晋书·河间平王颙传》，证据提示为“核心事项已有既有账本来源；本分解节点的参与者、执行范围和时间先后仍须逐案核验”。这些材料可支持年序、官方决定、战事或特定地点的线索，却难以直接统计所有人的财富、迁徙、信仰和动机。更稳妥的读法是区分诏令与实施、军报与人口、行政称谓与自我认同，并让地方材料的缺环限制解释的范围。

由此可见，该事件不是孤立的年表点，而是资源、信息与权力关系重新配置的一环。不同地区的官员、社群和家庭可能以不同方式承担征发、保护、迁移或宗教活动；现存来源只能照亮其中一部分，不能替沉默者制造统一的经验。

\eventanchor{JH-LEDGER-016-3}{永兴二年（305年）·东海王司马越以讨河间王司马颙为名起兵，八王之乱扩大为跨州郡的长期战争}账本记录西晋在永兴二年（305年）发生“东海王司马越以讨河间王司马颙为名起兵，八王之乱扩大为跨州郡的长期战争”。条目把背景概括为“东海王司马越以讨河间王司马颙为名起兵，八王之乱扩大为跨州郡的长期战争”发生的政治与区域条件共同作用，并指出该节点改变后续资源、联盟或制度处置的条件。这个节点进入区域社会时，要经由文书、道路、军队、市场、寺院和家庭等不同中介；因此，政治行动的宣布、地方接收和日常生活的改变不能被视作同一时刻完成。

本条依据《晋书·惠帝纪》；《晋书·东海孝献王越传》；《晋书·河间平王颙传》，证据提示为“核心事项已有既有账本来源；本分解节点的参与者、执行范围和时间先后仍须逐案核验”。这些材料可支持年序、官方决定、战事或特定地点的线索，却难以直接统计所有人的财富、迁徙、信仰和动机。更稳妥的读法是区分诏令与实施、军报与人口、行政称谓与自我认同，并让地方材料的缺环限制解释的范围。

由此可见，该事件不是孤立的年表点，而是资源、信息与权力关系重新配置的一环。不同地区的官员、社群和家庭可能以不同方式承担征发、保护、迁移或宗教活动；现存来源只能照亮其中一部分，不能替沉默者制造统一的经验。

\eventanchor{JH-LEDGER-016-4}{永兴二年（305年）·“东海王司马越以讨河间王司马颙为名起兵，八王之乱扩大为跨州郡的长期战争”的执行中的空间与人员处置}账本记录西晋在永兴二年（305年）发生““东海王司马越以讨河间王司马颙为名起兵，八王之乱扩大为跨州郡的长期战争”的执行中的空间与人员处置”。条目把背景概括为“东海王司马越以讨河间王司马颙为名起兵，八王之乱扩大为跨州郡的长期战争”发生的政治与区域条件共同作用，并指出该节点改变后续资源、联盟或制度处置的条件。这个节点进入区域社会时，要经由文书、道路、军队、市场、寺院和家庭等不同中介；因此，政治行动的宣布、地方接收和日常生活的改变不能被视作同一时刻完成。

本条依据《晋书·惠帝纪》；《晋书·东海孝献王越传》；《晋书·河间平王颙传》，证据提示为“核心事项已有既有账本来源；本分解节点的参与者、执行范围和时间先后仍须逐案核验”。这些材料可支持年序、官方决定、战事或特定地点的线索，却难以直接统计所有人的财富、迁徙、信仰和动机。更稳妥的读法是区分诏令与实施、军报与人口、行政称谓与自我认同，并让地方材料的缺环限制解释的范围。

由此可见，该事件不是孤立的年表点，而是资源、信息与权力关系重新配置的一环。不同地区的官员、社群和家庭可能以不同方式承担征发、保护、迁移或宗教活动；现存来源只能照亮其中一部分，不能替沉默者制造统一的经验。

\eventanchor{JH-LEDGER-016-5}{永兴二年（305年）·“东海王司马越以讨河间王司马颙为名起兵，八王之乱扩大为跨州郡的长期战争”的后续制度或军政调整}账本记录西晋在永兴二年（305年）发生““东海王司马越以讨河间王司马颙为名起兵，八王之乱扩大为跨州郡的长期战争”的后续制度或军政调整”。条目把背景概括为核心节点发生后仍须处理其遗留问题，并指出后续影响须在不同地区和材料类型中分别核验。这个节点进入区域社会时，要经由文书、道路、军队、市场、寺院和家庭等不同中介；因此，政治行动的宣布、地方接收和日常生活的改变不能被视作同一时刻完成。

本条依据《晋书·惠帝纪》；《晋书·东海孝献王越传》；《晋书·河间平王颙传》，证据提示为“核心事项已有既有账本来源；本分解节点的参与者、执行范围和时间先后仍须逐案核验”。这些材料可支持年序、官方决定、战事或特定地点的线索，却难以直接统计所有人的财富、迁徙、信仰和动机。更稳妥的读法是区分诏令与实施、军报与人口、行政称谓与自我认同，并让地方材料的缺环限制解释的范围。

由此可见，该事件不是孤立的年表点，而是资源、信息与权力关系重新配置的一环。不同地区的官员、社群和家庭可能以不同方式承担征发、保护、迁移或宗教活动；现存来源只能照亮其中一部分，不能替沉默者制造统一的经验。

\eventanchor{JH-LEDGER-016-6}{永兴二年（305年）·“东海王司马越以讨河间王司马颙为名起兵，八王之乱扩大为跨州郡的长期战争”的异源材料与影响范围的复核}账本记录西晋在永兴二年（305年）发生““东海王司马越以讨河间王司马颙为名起兵，八王之乱扩大为跨州郡的长期战争”的异源材料与影响范围的复核”。条目把背景概括为核心节点发生后仍须处理其遗留问题，并指出后续影响须在不同地区和材料类型中分别核验。这个节点进入区域社会时，要经由文书、道路、军队、市场、寺院和家庭等不同中介；因此，政治行动的宣布、地方接收和日常生活的改变不能被视作同一时刻完成。

本条依据《晋书·惠帝纪》；《晋书·东海孝献王越传》；《晋书·河间平王颙传》，证据提示为“核心事项已有既有账本来源；本分解节点的参与者、执行范围和时间先后仍须逐案核验”。这些材料可支持年序、官方决定、战事或特定地点的线索，却难以直接统计所有人的财富、迁徙、信仰和动机。更稳妥的读法是区分诏令与实施、军报与人口、行政称谓与自我认同，并让地方材料的缺环限制解释的范围。

由此可见，该事件不是孤立的年表点，而是资源、信息与权力关系重新配置的一环。不同地区的官员、社群和家庭可能以不同方式承担征发、保护、迁移或宗教活动；现存来源只能照亮其中一部分，不能替沉默者制造统一的经验。

\eventanchor{JH-LEDGER-017-1}{光熙元年（306年）·“惠帝在洛阳去世，司马炽即位为怀帝，司马越继续掌握主要军政权”的前期动员与资源准备}账本记录西晋在光熙元年（306年）发生““惠帝在洛阳去世，司马炽即位为怀帝，司马越继续掌握主要军政权”的前期动员与资源准备”。条目把背景概括为当时的权力、资源与信息约束推动相关过程展开，并指出这一过程为“惠帝在洛阳去世，司马炽即位为怀帝，司马越继续掌握主要军政权”提供可观察的条件。这个节点进入区域社会时，要经由文书、道路、军队、市场、寺院和家庭等不同中介；因此，政治行动的宣布、地方接收和日常生活的改变不能被视作同一时刻完成。

本条依据《晋书·惠帝纪》；《晋书·怀帝纪》；《晋书·东海孝献王越传》，证据提示为“核心事项已有既有账本来源；本分解节点的参与者、执行范围和时间先后仍须逐案核验”。这些材料可支持年序、官方决定、战事或特定地点的线索，却难以直接统计所有人的财富、迁徙、信仰和动机。更稳妥的读法是区分诏令与实施、军报与人口、行政称谓与自我认同，并让地方材料的缺环限制解释的范围。

由此可见，该事件不是孤立的年表点，而是资源、信息与权力关系重新配置的一环。不同地区的官员、社群和家庭可能以不同方式承担征发、保护、迁移或宗教活动；现存来源只能照亮其中一部分，不能替沉默者制造统一的经验。

\eventanchor{JH-LEDGER-017-2}{光熙元年（306年）·“惠帝在洛阳去世，司马炽即位为怀帝，司马越继续掌握主要军政权”的命令、联盟或地方响应}账本记录西晋在光熙元年（306年）发生““惠帝在洛阳去世，司马炽即位为怀帝，司马越继续掌握主要军政权”的命令、联盟或地方响应”。条目把背景概括为当时的权力、资源与信息约束推动相关过程展开，并指出这一过程为“惠帝在洛阳去世，司马炽即位为怀帝，司马越继续掌握主要军政权”提供可观察的条件。这个节点进入区域社会时，要经由文书、道路、军队、市场、寺院和家庭等不同中介；因此，政治行动的宣布、地方接收和日常生活的改变不能被视作同一时刻完成。

本条依据《晋书·惠帝纪》；《晋书·怀帝纪》；《晋书·东海孝献王越传》，证据提示为“核心事项已有既有账本来源；本分解节点的参与者、执行范围和时间先后仍须逐案核验”。这些材料可支持年序、官方决定、战事或特定地点的线索，却难以直接统计所有人的财富、迁徙、信仰和动机。更稳妥的读法是区分诏令与实施、军报与人口、行政称谓与自我认同，并让地方材料的缺环限制解释的范围。

由此可见，该事件不是孤立的年表点，而是资源、信息与权力关系重新配置的一环。不同地区的官员、社群和家庭可能以不同方式承担征发、保护、迁移或宗教活动；现存来源只能照亮其中一部分，不能替沉默者制造统一的经验。

\eventanchor{JH-LEDGER-017-3}{光熙元年（306年）·惠帝在洛阳去世，司马炽即位为怀帝，司马越继续掌握主要军政权}账本记录西晋在光熙元年（306年）发生“惠帝在洛阳去世，司马炽即位为怀帝，司马越继续掌握主要军政权”。条目把背景概括为“惠帝在洛阳去世，司马炽即位为怀帝，司马越继续掌握主要军政权”发生的政治与区域条件共同作用，并指出该节点改变后续资源、联盟或制度处置的条件。这个节点进入区域社会时，要经由文书、道路、军队、市场、寺院和家庭等不同中介；因此，政治行动的宣布、地方接收和日常生活的改变不能被视作同一时刻完成。

本条依据《晋书·惠帝纪》；《晋书·怀帝纪》；《晋书·东海孝献王越传》，证据提示为“核心事项已有既有账本来源；本分解节点的参与者、执行范围和时间先后仍须逐案核验”。这些材料可支持年序、官方决定、战事或特定地点的线索，却难以直接统计所有人的财富、迁徙、信仰和动机。更稳妥的读法是区分诏令与实施、军报与人口、行政称谓与自我认同，并让地方材料的缺环限制解释的范围。

由此可见，该事件不是孤立的年表点，而是资源、信息与权力关系重新配置的一环。不同地区的官员、社群和家庭可能以不同方式承担征发、保护、迁移或宗教活动；现存来源只能照亮其中一部分，不能替沉默者制造统一的经验。

\eventanchor{JH-LEDGER-017-4}{光熙元年（306年）·“惠帝在洛阳去世，司马炽即位为怀帝，司马越继续掌握主要军政权”的执行中的空间与人员处置}账本记录西晋在光熙元年（306年）发生““惠帝在洛阳去世，司马炽即位为怀帝，司马越继续掌握主要军政权”的执行中的空间与人员处置”。条目把背景概括为“惠帝在洛阳去世，司马炽即位为怀帝，司马越继续掌握主要军政权”发生的政治与区域条件共同作用，并指出该节点改变后续资源、联盟或制度处置的条件。这个节点进入区域社会时，要经由文书、道路、军队、市场、寺院和家庭等不同中介；因此，政治行动的宣布、地方接收和日常生活的改变不能被视作同一时刻完成。

本条依据《晋书·惠帝纪》；《晋书·怀帝纪》；《晋书·东海孝献王越传》，证据提示为“核心事项已有既有账本来源；本分解节点的参与者、执行范围和时间先后仍须逐案核验”。这些材料可支持年序、官方决定、战事或特定地点的线索，却难以直接统计所有人的财富、迁徙、信仰和动机。更稳妥的读法是区分诏令与实施、军报与人口、行政称谓与自我认同，并让地方材料的缺环限制解释的范围。

由此可见，该事件不是孤立的年表点，而是资源、信息与权力关系重新配置的一环。不同地区的官员、社群和家庭可能以不同方式承担征发、保护、迁移或宗教活动；现存来源只能照亮其中一部分，不能替沉默者制造统一的经验。

\eventanchor{JH-LEDGER-017-5}{光熙元年（306年）·“惠帝在洛阳去世，司马炽即位为怀帝，司马越继续掌握主要军政权”的后续制度或军政调整}账本记录西晋在光熙元年（306年）发生““惠帝在洛阳去世，司马炽即位为怀帝，司马越继续掌握主要军政权”的后续制度或军政调整”。条目把背景概括为核心节点发生后仍须处理其遗留问题，并指出后续影响须在不同地区和材料类型中分别核验。这个节点进入区域社会时，要经由文书、道路、军队、市场、寺院和家庭等不同中介；因此，政治行动的宣布、地方接收和日常生活的改变不能被视作同一时刻完成。

本条依据《晋书·惠帝纪》；《晋书·怀帝纪》；《晋书·东海孝献王越传》，证据提示为“核心事项已有既有账本来源；本分解节点的参与者、执行范围和时间先后仍须逐案核验”。这些材料可支持年序、官方决定、战事或特定地点的线索，却难以直接统计所有人的财富、迁徙、信仰和动机。更稳妥的读法是区分诏令与实施、军报与人口、行政称谓与自我认同，并让地方材料的缺环限制解释的范围。

由此可见，该事件不是孤立的年表点，而是资源、信息与权力关系重新配置的一环。不同地区的官员、社群和家庭可能以不同方式承担征发、保护、迁移或宗教活动；现存来源只能照亮其中一部分，不能替沉默者制造统一的经验。

\eventanchor{JH-LEDGER-017-6}{光熙元年（306年）·“惠帝在洛阳去世，司马炽即位为怀帝，司马越继续掌握主要军政权”的异源材料与影响范围的复核}账本记录西晋在光熙元年（306年）发生““惠帝在洛阳去世，司马炽即位为怀帝，司马越继续掌握主要军政权”的异源材料与影响范围的复核”。条目把背景概括为核心节点发生后仍须处理其遗留问题，并指出后续影响须在不同地区和材料类型中分别核验。这个节点进入区域社会时，要经由文书、道路、军队、市场、寺院和家庭等不同中介；因此，政治行动的宣布、地方接收和日常生活的改变不能被视作同一时刻完成。

本条依据《晋书·惠帝纪》；《晋书·怀帝纪》；《晋书·东海孝献王越传》，证据提示为“核心事项已有既有账本来源；本分解节点的参与者、执行范围和时间先后仍须逐案核验”。这些材料可支持年序、官方决定、战事或特定地点的线索，却难以直接统计所有人的财富、迁徙、信仰和动机。更稳妥的读法是区分诏令与实施、军报与人口、行政称谓与自我认同，并让地方材料的缺环限制解释的范围。

由此可见，该事件不是孤立的年表点，而是资源、信息与权力关系重新配置的一环。不同地区的官员、社群和家庭可能以不同方式承担征发、保护、迁移或宗教活动；现存来源只能照亮其中一部分，不能替沉默者制造统一的经验。
